% Included by manuscript.tex via % Included by manuscript.tex via % Included by manuscript.tex via % Included by manuscript.tex via \input{validation_section}.
% Not standalone: no preamble, no \begin{document}.
% Every number here was measured against an INDEPENDENT implementation
% or an analytic result -- never against a previous version of this code.

\section{Validation}
\label{sec:validation}

The implementation was checked against analytic solutions, against published
tabulations, and against four independently written software packages. We
report the comparisons in full, including the two cases in which they
uncovered defects in our own code, because the pattern is itself informative:
every defect found late in this work was invisible to internal consistency
checks and was exposed only by an external reference.

\subsection{Analytic and literature references}

\begin{table}[htbp]
\centering
\caption{Comparison with analytic results and published tabulations. Where a
difference is quoted it is the deviation from the reference, not from a
previous version of this code.}
\label{tab:analytic}
\small
\begin{tabular}{llll}
\toprule
Reference & Quantity & Ours & Agreement\\
\midrule
\citet{Wertheim1963}, \citet{Thiele1963} & $c(r)$, PY hard spheres & numerical & first order in $\Delta r$\\
                               & $g(\sigma^+)$, $\eta=0.1$--$0.4$ & numerical & $0.1\,\%$\\
                               & $S(0)=(1-\eta)^4/(1+2\eta)^2$ & numerical & converges\\
Analytic PY, $\eta=0.3$        & $\chi^{-1}$ & 11.07 & 10.66 (extrapolation)\\
                               & $\beta P/\rho$ & 3.890 & 3.816 (grid)\\
                               & PY inconsistency & $+1.33$ & $+1.25$\\
D'Aguanno \& Klein (1991)      & $1/S(0)$, excess pressure & --- & $0.01\,\%$, $0.02\,\%$\\
                               & excess energy, $g_{\max}$ & --- & $0.27\,\%$, $0.54\,\%$\\
\citet{Vrij1979} mixture PY         & $S^{\mathrm{AL}}_{ij}$, 3 classes & numerical & $0.09\,\%$\\
\bottomrule
\end{tabular}
\end{table}

\subsection{Comparison with independent software}

\begin{table}[htbp]
\centering
\caption{Comparison with independently written packages. ``exact'' denotes
agreement to all printed digits; the remaining entries are limited by
discretisation and converge as $\Delta r$ is reduced.}
\label{tab:software}
\small
\begin{tabular}{lllc}
\toprule
Package & Model & Method & Agreement\\
\midrule
\code{mixscatter} \citep{DiazMaier2024} & mixture PY \citep{Vrij1979} & analytic & $0.09\,\%$\\
\code{jscatter} & one-component PY & analytic & $0.02$--$0.12\,\%$\\
\code{sasmodels}, \code{jscatter} & RMSA \citep{HayterPenfold1981} & analytic & $10^{-12}$ (6/6)\\
\code{sasmodels} & two-Yukawa MSA \citep{Liu2005} & analytic & $0.21\,\%$\\
\code{jscatter} & adhesive hard spheres & analytic & see \S\ref{sec:welltrap}\\
\code{jscatter} & OZ solver, PY and HNC & numerical & $0.2$--$0.4\,\%$\\
\code{OrnsteinZernike.jl} & 4 bridge functions & analytic & $\sim10^{-16}$\\
\bottomrule
\end{tabular}
\end{table}

\paragraph{An absolute check: the compressibility.} Every comparison above
sets one calculation against another. One quantity admits an exact answer
instead. For monodisperse hard spheres the Percus--Yevick compressibility is
\begin{equation}
  S(0) = \frac{(1-\varphi)^{4}}{(1+2\varphi)^{2}},
  \label{eq:pyS0}
\end{equation}
so extrapolating the computed partials to the origin and comparing measures
absolute accuracy rather than agreement.

The result is that the transform type matters far more than its place in a
cost table suggests. At $\varphi = 0.40$ and 100 radial points per diameter,
the first-order DST-I gives $S(0)$ low by $5.27\,\%$; the second-order DST-IV
is low by $0.036\,\%$ on the same grid, and by $0.004\,\%$ at four times the
resolution. Type 1 refines as $3.96$ per fourfold step, textbook first order,
so it reaches $1.33\,\%$ where type 4 began 150 times better.

The cause is where the discontinuity lands. A DST-I places nodes at
$r = (n+1)\Delta r$, so the hard core at $r = \sigma$ falls ON a node and is
carried by a point that is neither inside nor outside; the excluded volume is
then wrong at first order in $\Delta r$. A DST-IV places them at
$(n+\frac{1}{2})\Delta r$, the core falls between points, and the step is
resolved to second order.

Two things make this worth reporting rather than absorbing into a
recommendation. First, the error is independent of the real-space range:
extending $r_{\max}$ from $41\sigma$ to $655\sigma$ left $S(0)$ unchanged in
the sixth decimal. A convergence study that refines resolution sees orderly
first-order behaviour and concludes the method works --- which it does, slowly,
towards an answer whose leading error is set by grid alignment. Second,
$S(0)$ is the isothermal compressibility, so this is a systematic bias in
exactly the quantity the volume fraction controls, and a fit will move
$\varphi$ to absorb it.

It also accounts for a discrepancy that had resisted explanation: a
multicomponent solve converged to a fixed point at $7\times10^{-12}$ yet
differing from two independent analytic references by $\sim10^{-3}$, the
difference immune to refinement and growing with $\varphi$. That is the same
effect seen from the other side.

The practical consequence is a restriction rather than a fix. Type 4 needs
every pair core between grid points, which one size class can be given and a
mixture cannot: the pair diameters $\sigma_{ij} = (\sigma_i+\sigma_j)/2$ come
from a moment-matched quadrature and are mutually incommensurate. The
polydisperse route is therefore first order, and until a mixture-capable
second-order treatment exists a dense polydisperse fit should use the finest
affordable resolution, with its fitted volume fraction read accordingly.

\paragraph{Reproducibility of the figures and of the table.} Three gaps are
recorded here rather than left to be discovered.

The remaining \code{jscatter} rows above --- one-component PY, adhesive
hard spheres and the numerical OZ comparison --- were measured where
\code{jscatter} was installed. Its dependency chain does not build under
MinGW, so they cannot be re-run on the machine this work is written on. Two
rows have since been brought back within reach, and by the same move:
replacing a reference that has to be obtained with one that can be
installed.

\paragraph{Two reference solvers replaced by \code{sasmodels}.} The RMSA
and two-Yukawa rows are now measured against SasView's \code{sasmodels},
which is 3-clause BSD and installs with \code{pip} alone. Both previously
depended on something a reader could not simply obtain --- \code{jscatter},
which will not build here, and a 97\,kB vendored copy of the two-Yukawa
solver used under a citation-ware grant the authors gave in writing. A
licence governs an implementation rather than a method, so \citet{Liu2005}
is cited exactly as before; what changes is that anyone can now repeat the
comparison.

Agreement is to $2\times10^{-8}$, $7\times10^{-8}$ and $3\times10^{-8}$ on
the three two-Yukawa state points where both solve cleanly. Of the
remaining two, \code{sasmodels} returns a physical $S(q)$ at
$\varphi = 0.10$ where the vendored solver reports no root found --- a
root-finder limitation rather than a statement about the physics, since the
same parameters at HIGHER volume fraction solve in both. The last,
$Z_1 = 6$ against $Z_2 = 4$, agrees only to $2\times10^{-3}$, with the
discrepancy peaking at the structure-factor maximum; the two inverse ranges
are close there, which is where the quartic's roots crowd together and two
implementations may select marginally differently. Both results are
physical and both are defensible.

\paragraph{The sign of $K$, and where the inconsistency actually was.}
\citet{Liu2005} write the tail as $-K\exp[-Z(r-1)]/r$, so positive $K$ is
ATTRACTIVE. This package's solver follows that convention, and so does
\code{sasmodels}. The vendored reference file did not: reproducing its
output required flipping both signs, after which the two agree to
$10^{-8}$, where every other candidate (length unit, pair ordering, $Z$
scaling) stayed above 20 per cent at every state point.

So the inconsistency was in the comparison apparatus rather than in the
solver being compared --- which is the more flattering of the two
possibilities, and was not the one assumed while the discrepancy was being
chased. Retiring the vendored file removes it entirely, and is a further
reason to prefer \code{sasmodels}: one fewer convention in the tree.

\code{fig3\_oscillation.pdf} has no generator in the repository. The figure
exists and nothing can reproduce it, which means a referee asking for any
change to it would find there is nothing to change.
\code{fig\_survey.pdf} names \code{tools/section53\_survey.py} in its own
caption and is the model the others should follow;
\code{fig1\_discretisation.pdf} now has
\code{tools/fig1\_discretisation.py}.

Regenerating it also caught a discrepancy in the figure itself.
\code{polydisperse\_nodes} offers two discretisations --- \code{sizeClasses},
the moment-matched generalised Gauss--Laguerre rule the caption describes,
and \code{quantileClasses}, which is quantile-based --- and only the first
reproduces the moments exactly. At $p = 3$ the Gauss--Laguerre rule returns
$\langle\sigma\rangle = 1.000000$ and $\langle\sigma^{2}\rangle = 1.062500$,
exactly $1 + s^{2}$; the quantile rule returns $0.9794$ and $0.9596$. Both
are usable discretisations, but only the first supports the claim printed
beside the figure. The moments are now shown on each panel so that the claim
can be checked rather than taken on trust.

\paragraph{The table's figures depend on the transform, and the table does
not say so.} The one-component Percus--Yevick row is quoted as
$0.02$--$0.12\,\%$ over $0.2 \le \varphi \le 0.4$, and that is a DST-IV
result. With the first-order DST-I at the same 100 radial points per
diameter the same comparison gives $0.70$, $1.37$ and $2.33\,\%$ at
$\varphi = 0.2$, $0.3$ and $0.4$ --- an order of magnitude worse, rising with
density, and reducible only fourfold per fourfold refinement
($0.17$, $0.33$, $0.51\,\%$ at 400 points).

This is the grid-alignment effect described above --- the hard core falling ON
a DST-I node rather than between DST-IV nodes --- seen in a second
quantity, and it carries the same caveat: type 4 is available for ONE size
class and not for a mixture, so the polydisperse route --- which is what this
paper is about --- carries the first-order figures rather than the quoted
ones. A reader comparing packages should know which transform produced a
number before comparing it with their own.

\paragraph{RMSA, and a units error that imitated a physics disagreement.}
\label{sec:rmsarescale}
This row agrees with \emph{two} independent implementations to $10^{-12}$
over $\gamma = 3.7$ to $816$ and $\kappa\sigma = 1.5$ to $3.8$, including the
strongly coupled points at which the Hayter--Penfold rescaling engages.
Reaching that took an afternoon, and the detour is worth recording because
the symptom was convincingly misleading.

The earlier claim was ``exact (9/9)'' against \code{jscatter}, which does not
build under MinGW: unrepeatable on the machine this is developed on.
\code{sasmodels}' \code{hayter\_msa} is a better reference --- a different
code lineage, installable with \code{pip} and nothing else --- but comparing
against it gave disagreements of 6 to 32 per cent that GREW WITH COUPLING.
That is the signature of a physics difference rather than a units slip, and
it was pursued as one: root selection between the two solutions the closure
admits, the rescaling threshold, radius-versus-diameter in $\gamma$. Each
was plausible, each was testable, each was wrong. No single rescaling of any
input reconciled the two, and the best-fitting factor itself drifted with
coupling.

The answer is in Eq.~(3a) of \citet{HayterPenfold1981}. The reduced pair
potential is written
\begin{equation}
\beta U(x) = \gamma\,\frac{e^{-kx}}{x},
\qquad x = r/\sigma,\quad k = \kappa\sigma,
\label{eq:hpgamma}
\end{equation}
so the value AT CONTACT is $\gamma e^{-k}$, and that --- not $\gamma$ --- is
what the paper equates to $\beta z^2/[\pi\epsilon_0\epsilon\sigma(2+k)^2]$.
\code{sasmodels} computes the contact potential; \code{jscatter} and the
wrapper here take the amplitude $\gamma$. The two differ by $e^{k}$, and
because $k$ varies with the screening the discrepancy grows with coupling
--- which is exactly why it read as physics.

Two things are worth taking from this. A discrepancy that scales with a
physical parameter is NOT evidence that the difference is physical: a
conversion factor depending on that parameter behaves identically. And the
arbiter was the original paper rather than either code's documentation ---
both describe $\gamma$ as ``the contact potential'', and both use the phrase
loosely for a quantity their own sources define differently.

With the factor applied, \code{tools/rmsa\_vs\_sasmodels.py} reproduces
\code{sasmodels} at all six state points, and the same inputs reproduce
\code{jscatter} to the same precision. The claim in the table is therefore
checkable by any reader with \code{pip install sasmodels}, which the one it
replaces was not.

\paragraph{A closure whose domain is unmapped.} The Carbajal--Tinoco bridge
function \citep{CarbajalTinoco2022} is implicit, $b = T(\gamma + b)$, and
has no solution at all below a
threshold: substituting $w = \gamma + b$ turns it into the explicit map
$\gamma(w) = w - T(w)$, whose range is bounded below. Numerically
$\gamma_{\min} = -0.509$ at $\lambda = 0$, rising towards zero as $\lambda$
grows, with two solutions above it and none below. Neither the source paper
nor the comparison of \citet{PihlajamaaJanssen2024} locates that boundary;
the latter simply excludes closures from its figures ``due to convergence
issues at this state point''.

\subsection{Analytic references implemented in this work}
\label{sec:analyticrefs}

Two of the three classical solvable polydisperse models had no reference
implementation available to us, so we implemented them. Together with the
Vrij hard-sphere solution supplied by \code{mixscatter}, they give exact
analytic coverage of the polydisperse hard-sphere, charged-hard-sphere and
sticky-hard-sphere fluids, against which the numerical solver can be checked
at any state point.

\paragraph{Polydisperse charged hard spheres.}
Appendix A of Gazzillo, Giacometti \& Carsughi (1999), the corrected form of
their 1997 paper --- the correction supplies a factor $\pi$ omitted from
$P_z$ and repairs an error in the determinant. This is the same physics
implemented by the FLAC program of Carsughi \emph{et al.}, which cites the
uncorrected 1997 paper. Our implementation reproduces the neutral
monodisperse limit against analytic Percus--Yevick to $3\times10^{-14}$, the
neutral polydisperse limit against \code{mixscatter} to $1\times10^{-13}$,
and, for charged systems, all three effects the authors report for increasing
polydispersity: the first peak is lowered ($1.7350\to1.3546$), shifted to
smaller $Q$ ($5.841\to5.030$ in units of $Q\langle\sigma\rangle$) and
$S(0)$ increased ($0.0187\to0.0484$), reproducing their component counts (79,
112, 149) exactly.

One limitation of that route is worth stating, because it is a modelling
assumption rather than an approximation that improves with effort. The
corresponding-states construction takes the valence to be PROPORTIONAL TO
SURFACE AREA, so charge polydispersity is fully correlated with size
polydispersity: one distribution governs both. A sample whose surface
chemistry does not scale with area --- differing grafting densities across a
size fraction, say --- cannot be represented.

\citet{GinozaYasutomi1998} treat size and interaction polydispersity as
\emph{independent}, on the same factorisable-coefficient mean-spherical
solution, and would be the route to take if that case ever arose. Both
descend from \citet{Ginoza1986}, who showed that factorisable coefficients
$K_{ij} = K d_i d_j$ collapse the multicomponent Blum--H\o ye equations to
rational functions of a single scalar $\Gamma$ --- which is what makes any of
this affordable for a size distribution, since each component then enters
only through $(\sigma_i, \rho_i, d_i)$ and the cost per evaluation is linear
in the number of classes rather than quadratic.

A partial implementation of the Ginoza route exists in
\code{ginoza\_yukawa\_msa.py} and is deliberately shelved: it solves for
$\Gamma$ and stops short of $S(q)$. Finishing it would reproduce, from
equations transcribed out of a scanned 1986 paper, a structure factor the
Gazzillo route already computes and which is validated above at machine
precision. The two closing conditions Ginoza gives, his Eqs.~(38a) and
(38b), are both implemented there and agree to $10^{-12}$, so the
transcription is at least self-consistent; that is recorded in case the
decoupled-charge case ever makes the module worth completing.

\paragraph{Polydisperse sticky hard spheres.}
\citet{GazzilloGiacometti2000}, a mean-spherical solution in the Baxter sticky
limit with factorisable coupling parameters $K_{ij}=Y_iY_j$, which renders
the Baxter factor matrix three-dyadic and collapses the $p\times p$ inversion
to a $4\times4$ determinant for any number of components. Validation:
the $\gamma_0=0$ limits against analytic Percus--Yevick ($4\times10^{-14}$)
and \code{mixscatter} ($1\times10^{-13}$); the monodisperse sticky
compressibility against their closed form ($3\times10^{-8}$); the critical
point $\eta_c=(\sqrt3-1)/2$, $\gamma_{0c}^2=(\sqrt3+2)/6$ located as the
divergence of $S(0)$; and the generalised Boyle temperatures of their
Eq.~(53) to the printed digits ($2.787$, $1.677$, $2.986$, $2.897$ against
$2.79$, $1.68$, $2.99$, $2.90$).

Note that the stickiness parameter of this model is \emph{not} Baxter's
$\tau$ \citep{Baxter1968}: Baxter's Percus--Yevick solution for mixtures requires $p(p+1)/2$
coupled quadratic equations and is not in general dyadic, which is why the
mean-spherical model exists. The two agree only in the hard-sphere limit.

Three entries in Table~\ref{tab:software} deserve comment.

The \textbf{RMSA} comparison was exact at all nine state points tested
($\eta = 0.05$--$0.25$, screening lengths $2$--$10$\,nm, contact potentials
$20$--$100\,k_BT$), including four at which the Hayter--Penfold rescaling
\citep{HayterPenfold1981}
procedure engaged, so the agreement is not merely that of the underlying MSA
branch. Two independently written implementations of a numerically delicate
algorithm reproduce each other to all printed digits.

The \textbf{two-Yukawa} comparison validates the hard-core double-Yukawa
potential, including cases of mixed sign (attractive $K_1$, repulsive $K_2$),
which is the regime in which a sign convention is most easily mistaken. The
reference implementation states the identical convention,
$V(r) = -k_BT[K_1 e^{-Z_1(r-1)} + K_2 e^{-Z_2(r-1)}]/r$ for $r>\sigma$.

The \textbf{numerical} comparison is the only external check available for the
hypernetted-chain closure, which has no analytic solution. Both solvers are
discretised and differ in transform, mixing scheme and hard-core
representation, so the $0.3\,\%$ agreement is the expected order.

\subsection{Two defects found by external comparison}

\paragraph{Identical hard cores in mixtures.}
Comparison with \code{mixscatter} revealed a discrepancy of $5.2\,\%$ that,
diagnostically, did \emph{not} decrease under grid refinement and was
\emph{linear in the size spread}, vanishing exactly in the monodisperse limit.
The cause was that the hard-sphere potential routine positioned its core by
array index rather than against the radial grid, so it did not see the
rescaled grid on which the polydisperse construction of
equation~(\ref{eq:reducedtail}) relies. Every pair therefore received its core
at the same radius: the mixture was being solved as though all particles were
identical. After correction the agreement is $0.09\,\%$ and converges. It was
the only one of twenty-one potentials written in that form.

\subsection{Two numerical traps in fitting}

Both of the following are properties of the fitting problem rather than of
this implementation, and both produce a converged fit with a plausible
$\chi^2$ and silently wrong parameters. We report them because neither is
evident from the literature and both cost us considerable time.

\paragraph{A rank-revealing linear solve can discard a parameter entirely.}
It is natural to split the fit: the scale, the flat background and any
power-law amplitude enter the intensity linearly and can be obtained exactly
by weighted least squares at every iteration, leaving only the shape
parameters to the nonlinear optimiser. The saving is real. The hazard is
that the design columns need not be comparable in magnitude. A model built
from scattering length densities in $\mathrm{cm}^{-2}$ carries the contrast
squared and so has magnitude $\sim10^{20}$, while the constant background
column is unity --- a condition number near $10^{20}$.
\code{numpy.linalg.lstsq}, like any rank-revealing solver, discards singular
values below $\max(M,N)\epsilon$, so the background and power-law
coefficients are returned as \emph{exactly} zero and the solve reports
success. The symptom is a fit that a hand-chosen background beats by tens of
percent in $\chi^2$, which is indistinguishable from a broken optimiser.
Equilibrating the columns to unit norm before the solve and rescaling
afterwards recovers all coefficients to five figures and changes nothing
where the problem was already well conditioned \citep{GolubVanLoan2013}.

The same scaling defeats the nonlinear optimiser: a parameter of order
$10^{10}$ satisfies \code{least\_squares}'s gradient test after a single
evaluation, so it never moves and is reported with a meaningless
uncertainty. Since the absolute scattering length densities are in any case
perfectly correlated with the scale, we fit their RATIO, which is of order
unity and whose sign is the quantity the data actually determine.

\paragraph{Gaussian resolution smearing fails where it is needed most.} A
small-angle instrument broadens $\mathbf{Q}$ in two dimensions in the
detector plane, and the data are then radially averaged; the distribution of
the magnitude about a true $Q_0$ is Rician, Eq.~\eqref{eq:rician} below, not
Gaussian \citep{Pedersen1990,BarkerPedersen1995}. A Gaussian in $Q$ is its
$Q_0/\sigma \gg 1$ limit and is entirely adequate wherever
$\mathrm{d}Q/Q$ is small.

\begin{equation}
  P(Q\,|\,Q_0) \;=\; \frac{Q}{\sigma^{2}}
  \exp\!\left(-\frac{Q^{2}+Q_0^{2}}{2\sigma^{2}}\right)
  I_0\!\left(\frac{QQ_0}{\sigma^{2}}\right)
  \label{eq:rician}
\end{equation}

It is not always small. For a spherical-shell dataset used to develop this
work, $\mathrm{d}Q/Q$ reaches $0.55$ at the lowest measured point, where a
Gaussian kernel places $44\,\%$ of that point's weight below $Q_{\min}$ and
some of it at \emph{negative} $Q$. The leading factor of $Q$ in
\eqref{eq:rician} suppresses the origin and excludes negative $Q$ by
construction, reducing that to $34\,\%$. Both forms reproduce the unsmeared
curve to $2\times10^{-16}$ as $\mathrm{d}Q \to 0$, which is the check that
catches a mis-normalised kernel. Evaluate \eqref{eq:rician} through the
exponentially scaled $\tilde{I}_0(x) = e^{-|x|}I_0(x)$: the argument reaches
several hundred at these widths, where $I_0$ overflows a double long before
the prefactor cancels it.

A related trap follows. A background term $A\,Q^{-4+d}$ diverges at the
origin while the smearing kernel reaches below the measured range, so the
smeared value at the lowest point is governed by whatever guard prevents
division by zero --- in our case $6.6\times10^{5}$ times the unsmeared value,
and varying by five orders of magnitude as that guard was moved. The power
law describes the measured range and says nothing about $Q \to 0$, where any
real curve turns over; holding it constant below the lowest measured $Q$
removes the arbitrariness.

\begin{quote}\small\itshape
Draft note on citations. Now placed: \code{OrnsteinZernike1914} and
\code{PercusYevick1958} (the equation and closure, \S2),
\code{HansenMcDonald2013}, \code{Anderson1965}, \code{WalkerNi2011},
\code{HayterPenfold1981}, \code{Liu2005}, \code{DiazMaier2024}
(\code{mixscatter}), \code{BarkerPedersen1995} and
\code{GolubVanLoan2013}.

Still in \code{ozgui.bib} and uncited: \code{DuhHenderson1996},
\code{Gazzillo1999}, \code{Gazzillo2006}, \code{KalyuzhnyiCummings2006},
\code{Carsughi2000}. Each has a natural home --- the first in the discussion
of mixing rules, the Gazzillo papers where stickiness and its size
dependence are treated --- but placing a citation where it is not actually
used is worse than leaving it out, so they are left for whoever writes those
sections.

Still absent and worth adding: software citations for NumPy, SciPy, SUNDIALS
and NLopt, all of which the work depends on.

One substantive point from the verification. \citet{BarkerPedersen1995}
conclude that a Gaussian approximation to the resolution is PREFERRED in most
experiments for ease of calculation, and is inadequate only in particular
cases --- of which they name Porod scattering. The argument above should be
framed as identifying one of those known exceptions rather than as a general
correction, which is both more defensible and more useful.
\end{quote}

\subsection{What these figures measure}

A note on what these figures measure. Every agreement quoted in this section
was obtained with the type-1 discrete sine transform, whose grid points fall
\emph{on} the hard core, and which is therefore first order in $\Delta r$. A
type-4 transform places the grid at $(n+\tfrac12)\Delta r$ so that the core
lies between points --- the midpoint rule --- and is second order. For one
component it is both more accurate and cheaper, and both advantages grow with
resolution:

\begin{center}\small
\begin{tabular}{rrrr}
\toprule
points per $\sigma$ & type-4 time & accuracy gain\\
\midrule
100 & $0.73\times$ & $139\times$\\
400 & $0.65\times$ & $547\times$\\
800 & $0.49\times$ & $1092\times$\\
\bottomrule
\end{tabular}
\end{center}

The times are relative to type 1 at the same points per diameter; the
accuracy gain is the ratio of the two errors. The growth of that ratio with
resolution is the second-order against first-order signature. The figures
quoted above would improve correspondingly.

The two transforms require different grid \emph{lengths}, which is easy to
overlook and costly to get wrong: a DST-I is computed through an FFT of
length $2(N+1)$ and so wants $N = 2^{k}-1$, whereas a DST-IV is a length-$N$
transform and wants $N = 2^{k}$. Sizing both for the former leaves the latter
on a composite length ($16383 = 3\times43\times127$) and costs it a factor of
three in the transform, enough to make the second-order method appear slower
than the first-order one. Accuracy is unaffected by the choice.

We nevertheless retain type 1 as the default, because the type-4 condition
cannot be met for a mixture. It requires every pair core $\sigma_{ij}$ to
fall between grid points, and with size classes drawn from a Gaussian
quadrature the $\sigma_i$ are irrational, so the $p(p+1)/2$ conditions cannot
be satisfied simultaneously by any choice of resolution. Type 4 is available
as an option and is correct for one component, where the single core
coincides with the reduced unit by construction.


The defect had been invisible to every existing test because the resulting
error vanishes in precisely the monodisperse limit where those tests looked.

\paragraph{Under-resolved narrow wells.}
\label{sec:welltrap}
Comparison with an analytic adhesive-sphere solution showed errors that are
governed not by the number of grid points per particle diameter but by the
number \emph{across the attractive well}:

\begin{center}\small
\begin{tabular}{cc}
\toprule
points across well & max.\ relative error\\
\midrule
2 & 0.916\\
4 & 0.324\\
8 & 0.142\\
16 & 0.068\\
32 & 0.035\\
\bottomrule
\end{tabular}
\end{center}

This is ordinary first-order convergence in $\Delta r/\delta$ and not a
defect, but the default discretisation violates it badly: at 100 points per
diameter a well of width $\delta = 0.02\sigma$ spans two grid points and the
structure factor is in error by ninety per cent, without any warning. We
recommend $\;\gtrsim 30\,\sigma/\delta$ points per diameter for any potential
with a narrow feature, and note that this requirement applies equally to
square-well and piecewise-constant potentials.

\subsection{Internal consistency}

The following are not independent checks but constrain the implementation:
reduction of the polydisperse construction to a single size class reproduces
the original one-component routines bit-identically; the extended
Rogers--Young closure reduces to Rogers--Young to $2.5\times10^{-16}$ when its
additional parameter vanishes; the hard-core double-Yukawa potential agrees
with the three-Yukawa form to $\sim10^{-16}$ under the stated
reparameterisation; and the exact intensity reproduces brute-force
integration of $\langle|F|^2\rangle$ in the dilute limit
(ripple $0.041$ against $0.038$).

% ---------------------------------------------------------------------
% References to add to the bibliography
% ---------------------------------------------------------------------
% Gazzillo, D., Giacometti, A. & Carsughi, F. (1997). J. Chem. Phys. 107,
%   10141. [polydisperse charged hard spheres, MSA]
% Gazzillo, D., Giacometti, A. & Carsughi, F. (1999). arXiv:cond-mat/9909153.
%   [corrected Appendix A; scaling vs decoupling approximation]
% Gazzillo, D. & Giacometti, A. (2000). arXiv:cond-mat/0011330.
%   [polydisperse sticky hard spheres, MSA, factorisable coupling]
% Carsughi, F., Giacometti, A. & Gazzillo, D. (2000). arXiv:cond-mat/0011242.
%   [the FLAC program]
% Tutschka, C. & Kahl, G. (1998). J. Chem. Phys. 108, 9498.
% Blum, L. & Hoye, J. S. (1978). J. Stat. Phys. 19, 317.
%   [multicomponent multi-Yukawa MSA; the finite-range generalisation]
% Hoye, J. S. & Stell, G. [SCOZA]; Reiner, A. & Hoye, J. S. (2008).
%   arXiv:0712.3739. [SCOZA near the critical point]
% Bergmann, M. J. et al. mixscatter (Python package).
%   [add version/DOI; used for the analytic Vrij mixture PY]
% Biehl, R. (2019). jscatter. PLoS ONE 14, e0218789.
%   [Two_Yukawa after Liu et al. (2005); RMSA after Hayter & Penfold (1981)]
% Hayter, J. B. & Penfold, J. (1981). Mol. Phys. 42, 109-118.
% Liu, Y., Chen, W.-R. & Chen, S.-H. (2005). J. Chem. Phys. 122, 044507.
%   [two-Yukawa analytic solution]
% Pihlajamaa, I. & Janssen, L. M. C. (2024). Phys. Rev. E 110, 044608.
%   [closure comparison; OrnsteinZernike.jl]
% Vrij, A. (1979). J. Chem. Phys. 71, 3267-3270.
%   [mixture PY, the reference for the mixture comparison]
% Wertheim, M. S. (1963). Phys. Rev. Lett. 10, 321.
% Thiele, E. (1963). J. Chem. Phys. 39, 474.
% Baxter, R. J. (1968). J. Chem. Phys. 49, 2770.
% Malijevsky, A. & Labik, S. (1987). Mol. Phys. 60, 663.
.
% Not standalone: no preamble, no \begin{document}.
% Every number here was measured against an INDEPENDENT implementation
% or an analytic result -- never against a previous version of this code.

\section{Validation}
\label{sec:validation}

The implementation was checked against analytic solutions, against published
tabulations, and against four independently written software packages. We
report the comparisons in full, including the two cases in which they
uncovered defects in our own code, because the pattern is itself informative:
every defect found late in this work was invisible to internal consistency
checks and was exposed only by an external reference.

\subsection{Analytic and literature references}

\begin{table}[htbp]
\centering
\caption{Comparison with analytic results and published tabulations. Where a
difference is quoted it is the deviation from the reference, not from a
previous version of this code.}
\label{tab:analytic}
\small
\begin{tabular}{llll}
\toprule
Reference & Quantity & Ours & Agreement\\
\midrule
\citet{Wertheim1963}, \citet{Thiele1963} & $c(r)$, PY hard spheres & numerical & first order in $\Delta r$\\
                               & $g(\sigma^+)$, $\eta=0.1$--$0.4$ & numerical & $0.1\,\%$\\
                               & $S(0)=(1-\eta)^4/(1+2\eta)^2$ & numerical & converges\\
Analytic PY, $\eta=0.3$        & $\chi^{-1}$ & 11.07 & 10.66 (extrapolation)\\
                               & $\beta P/\rho$ & 3.890 & 3.816 (grid)\\
                               & PY inconsistency & $+1.33$ & $+1.25$\\
D'Aguanno \& Klein (1991)      & $1/S(0)$, excess pressure & --- & $0.01\,\%$, $0.02\,\%$\\
                               & excess energy, $g_{\max}$ & --- & $0.27\,\%$, $0.54\,\%$\\
\citet{Vrij1979} mixture PY         & $S^{\mathrm{AL}}_{ij}$, 3 classes & numerical & $0.09\,\%$\\
\bottomrule
\end{tabular}
\end{table}

\subsection{Comparison with independent software}

\begin{table}[htbp]
\centering
\caption{Comparison with independently written packages. ``exact'' denotes
agreement to all printed digits; the remaining entries are limited by
discretisation and converge as $\Delta r$ is reduced.}
\label{tab:software}
\small
\begin{tabular}{lllc}
\toprule
Package & Model & Method & Agreement\\
\midrule
\code{mixscatter} \citep{DiazMaier2024} & mixture PY \citep{Vrij1979} & analytic & $0.09\,\%$\\
\code{jscatter} & one-component PY & analytic & $0.02$--$0.12\,\%$\\
\code{sasmodels}, \code{jscatter} & RMSA \citep{HayterPenfold1981} & analytic & $10^{-12}$ (6/6)\\
\code{sasmodels} & two-Yukawa MSA \citep{Liu2005} & analytic & $0.21\,\%$\\
\code{jscatter} & adhesive hard spheres & analytic & see \S\ref{sec:welltrap}\\
\code{jscatter} & OZ solver, PY and HNC & numerical & $0.2$--$0.4\,\%$\\
\code{OrnsteinZernike.jl} & 4 bridge functions & analytic & $\sim10^{-16}$\\
\bottomrule
\end{tabular}
\end{table}

\paragraph{An absolute check: the compressibility.} Every comparison above
sets one calculation against another. One quantity admits an exact answer
instead. For monodisperse hard spheres the Percus--Yevick compressibility is
\begin{equation}
  S(0) = \frac{(1-\varphi)^{4}}{(1+2\varphi)^{2}},
  \label{eq:pyS0}
\end{equation}
so extrapolating the computed partials to the origin and comparing measures
absolute accuracy rather than agreement.

The result is that the transform type matters far more than its place in a
cost table suggests. At $\varphi = 0.40$ and 100 radial points per diameter,
the first-order DST-I gives $S(0)$ low by $5.27\,\%$; the second-order DST-IV
is low by $0.036\,\%$ on the same grid, and by $0.004\,\%$ at four times the
resolution. Type 1 refines as $3.96$ per fourfold step, textbook first order,
so it reaches $1.33\,\%$ where type 4 began 150 times better.

The cause is where the discontinuity lands. A DST-I places nodes at
$r = (n+1)\Delta r$, so the hard core at $r = \sigma$ falls ON a node and is
carried by a point that is neither inside nor outside; the excluded volume is
then wrong at first order in $\Delta r$. A DST-IV places them at
$(n+\frac{1}{2})\Delta r$, the core falls between points, and the step is
resolved to second order.

Two things make this worth reporting rather than absorbing into a
recommendation. First, the error is independent of the real-space range:
extending $r_{\max}$ from $41\sigma$ to $655\sigma$ left $S(0)$ unchanged in
the sixth decimal. A convergence study that refines resolution sees orderly
first-order behaviour and concludes the method works --- which it does, slowly,
towards an answer whose leading error is set by grid alignment. Second,
$S(0)$ is the isothermal compressibility, so this is a systematic bias in
exactly the quantity the volume fraction controls, and a fit will move
$\varphi$ to absorb it.

It also accounts for a discrepancy that had resisted explanation: a
multicomponent solve converged to a fixed point at $7\times10^{-12}$ yet
differing from two independent analytic references by $\sim10^{-3}$, the
difference immune to refinement and growing with $\varphi$. That is the same
effect seen from the other side.

The practical consequence is a restriction rather than a fix. Type 4 needs
every pair core between grid points, which one size class can be given and a
mixture cannot: the pair diameters $\sigma_{ij} = (\sigma_i+\sigma_j)/2$ come
from a moment-matched quadrature and are mutually incommensurate. The
polydisperse route is therefore first order, and until a mixture-capable
second-order treatment exists a dense polydisperse fit should use the finest
affordable resolution, with its fitted volume fraction read accordingly.

\paragraph{Reproducibility of the figures and of the table.} Three gaps are
recorded here rather than left to be discovered.

The remaining \code{jscatter} rows above --- one-component PY, adhesive
hard spheres and the numerical OZ comparison --- were measured where
\code{jscatter} was installed. Its dependency chain does not build under
MinGW, so they cannot be re-run on the machine this work is written on. Two
rows have since been brought back within reach, and by the same move:
replacing a reference that has to be obtained with one that can be
installed.

\paragraph{Two reference solvers replaced by \code{sasmodels}.} The RMSA
and two-Yukawa rows are now measured against SasView's \code{sasmodels},
which is 3-clause BSD and installs with \code{pip} alone. Both previously
depended on something a reader could not simply obtain --- \code{jscatter},
which will not build here, and a 97\,kB vendored copy of the two-Yukawa
solver used under a citation-ware grant the authors gave in writing. A
licence governs an implementation rather than a method, so \citet{Liu2005}
is cited exactly as before; what changes is that anyone can now repeat the
comparison.

Agreement is to $2\times10^{-8}$, $7\times10^{-8}$ and $3\times10^{-8}$ on
the three two-Yukawa state points where both solve cleanly. Of the
remaining two, \code{sasmodels} returns a physical $S(q)$ at
$\varphi = 0.10$ where the vendored solver reports no root found --- a
root-finder limitation rather than a statement about the physics, since the
same parameters at HIGHER volume fraction solve in both. The last,
$Z_1 = 6$ against $Z_2 = 4$, agrees only to $2\times10^{-3}$, with the
discrepancy peaking at the structure-factor maximum; the two inverse ranges
are close there, which is where the quartic's roots crowd together and two
implementations may select marginally differently. Both results are
physical and both are defensible.

\paragraph{The sign of $K$, and where the inconsistency actually was.}
\citet{Liu2005} write the tail as $-K\exp[-Z(r-1)]/r$, so positive $K$ is
ATTRACTIVE. This package's solver follows that convention, and so does
\code{sasmodels}. The vendored reference file did not: reproducing its
output required flipping both signs, after which the two agree to
$10^{-8}$, where every other candidate (length unit, pair ordering, $Z$
scaling) stayed above 20 per cent at every state point.

So the inconsistency was in the comparison apparatus rather than in the
solver being compared --- which is the more flattering of the two
possibilities, and was not the one assumed while the discrepancy was being
chased. Retiring the vendored file removes it entirely, and is a further
reason to prefer \code{sasmodels}: one fewer convention in the tree.

\code{fig3\_oscillation.pdf} has no generator in the repository. The figure
exists and nothing can reproduce it, which means a referee asking for any
change to it would find there is nothing to change.
\code{fig\_survey.pdf} names \code{tools/section53\_survey.py} in its own
caption and is the model the others should follow;
\code{fig1\_discretisation.pdf} now has
\code{tools/fig1\_discretisation.py}.

Regenerating it also caught a discrepancy in the figure itself.
\code{polydisperse\_nodes} offers two discretisations --- \code{sizeClasses},
the moment-matched generalised Gauss--Laguerre rule the caption describes,
and \code{quantileClasses}, which is quantile-based --- and only the first
reproduces the moments exactly. At $p = 3$ the Gauss--Laguerre rule returns
$\langle\sigma\rangle = 1.000000$ and $\langle\sigma^{2}\rangle = 1.062500$,
exactly $1 + s^{2}$; the quantile rule returns $0.9794$ and $0.9596$. Both
are usable discretisations, but only the first supports the claim printed
beside the figure. The moments are now shown on each panel so that the claim
can be checked rather than taken on trust.

\paragraph{The table's figures depend on the transform, and the table does
not say so.} The one-component Percus--Yevick row is quoted as
$0.02$--$0.12\,\%$ over $0.2 \le \varphi \le 0.4$, and that is a DST-IV
result. With the first-order DST-I at the same 100 radial points per
diameter the same comparison gives $0.70$, $1.37$ and $2.33\,\%$ at
$\varphi = 0.2$, $0.3$ and $0.4$ --- an order of magnitude worse, rising with
density, and reducible only fourfold per fourfold refinement
($0.17$, $0.33$, $0.51\,\%$ at 400 points).

This is the grid-alignment effect described above --- the hard core falling ON
a DST-I node rather than between DST-IV nodes --- seen in a second
quantity, and it carries the same caveat: type 4 is available for ONE size
class and not for a mixture, so the polydisperse route --- which is what this
paper is about --- carries the first-order figures rather than the quoted
ones. A reader comparing packages should know which transform produced a
number before comparing it with their own.

\paragraph{RMSA, and a units error that imitated a physics disagreement.}
\label{sec:rmsarescale}
This row agrees with \emph{two} independent implementations to $10^{-12}$
over $\gamma = 3.7$ to $816$ and $\kappa\sigma = 1.5$ to $3.8$, including the
strongly coupled points at which the Hayter--Penfold rescaling engages.
Reaching that took an afternoon, and the detour is worth recording because
the symptom was convincingly misleading.

The earlier claim was ``exact (9/9)'' against \code{jscatter}, which does not
build under MinGW: unrepeatable on the machine this is developed on.
\code{sasmodels}' \code{hayter\_msa} is a better reference --- a different
code lineage, installable with \code{pip} and nothing else --- but comparing
against it gave disagreements of 6 to 32 per cent that GREW WITH COUPLING.
That is the signature of a physics difference rather than a units slip, and
it was pursued as one: root selection between the two solutions the closure
admits, the rescaling threshold, radius-versus-diameter in $\gamma$. Each
was plausible, each was testable, each was wrong. No single rescaling of any
input reconciled the two, and the best-fitting factor itself drifted with
coupling.

The answer is in Eq.~(3a) of \citet{HayterPenfold1981}. The reduced pair
potential is written
\begin{equation}
\beta U(x) = \gamma\,\frac{e^{-kx}}{x},
\qquad x = r/\sigma,\quad k = \kappa\sigma,
\label{eq:hpgamma}
\end{equation}
so the value AT CONTACT is $\gamma e^{-k}$, and that --- not $\gamma$ --- is
what the paper equates to $\beta z^2/[\pi\epsilon_0\epsilon\sigma(2+k)^2]$.
\code{sasmodels} computes the contact potential; \code{jscatter} and the
wrapper here take the amplitude $\gamma$. The two differ by $e^{k}$, and
because $k$ varies with the screening the discrepancy grows with coupling
--- which is exactly why it read as physics.

Two things are worth taking from this. A discrepancy that scales with a
physical parameter is NOT evidence that the difference is physical: a
conversion factor depending on that parameter behaves identically. And the
arbiter was the original paper rather than either code's documentation ---
both describe $\gamma$ as ``the contact potential'', and both use the phrase
loosely for a quantity their own sources define differently.

With the factor applied, \code{tools/rmsa\_vs\_sasmodels.py} reproduces
\code{sasmodels} at all six state points, and the same inputs reproduce
\code{jscatter} to the same precision. The claim in the table is therefore
checkable by any reader with \code{pip install sasmodels}, which the one it
replaces was not.

\paragraph{A closure whose domain is unmapped.} The Carbajal--Tinoco bridge
function \citep{CarbajalTinoco2022} is implicit, $b = T(\gamma + b)$, and
has no solution at all below a
threshold: substituting $w = \gamma + b$ turns it into the explicit map
$\gamma(w) = w - T(w)$, whose range is bounded below. Numerically
$\gamma_{\min} = -0.509$ at $\lambda = 0$, rising towards zero as $\lambda$
grows, with two solutions above it and none below. Neither the source paper
nor the comparison of \citet{PihlajamaaJanssen2024} locates that boundary;
the latter simply excludes closures from its figures ``due to convergence
issues at this state point''.

\subsection{Analytic references implemented in this work}
\label{sec:analyticrefs}

Two of the three classical solvable polydisperse models had no reference
implementation available to us, so we implemented them. Together with the
Vrij hard-sphere solution supplied by \code{mixscatter}, they give exact
analytic coverage of the polydisperse hard-sphere, charged-hard-sphere and
sticky-hard-sphere fluids, against which the numerical solver can be checked
at any state point.

\paragraph{Polydisperse charged hard spheres.}
Appendix A of Gazzillo, Giacometti \& Carsughi (1999), the corrected form of
their 1997 paper --- the correction supplies a factor $\pi$ omitted from
$P_z$ and repairs an error in the determinant. This is the same physics
implemented by the FLAC program of Carsughi \emph{et al.}, which cites the
uncorrected 1997 paper. Our implementation reproduces the neutral
monodisperse limit against analytic Percus--Yevick to $3\times10^{-14}$, the
neutral polydisperse limit against \code{mixscatter} to $1\times10^{-13}$,
and, for charged systems, all three effects the authors report for increasing
polydispersity: the first peak is lowered ($1.7350\to1.3546$), shifted to
smaller $Q$ ($5.841\to5.030$ in units of $Q\langle\sigma\rangle$) and
$S(0)$ increased ($0.0187\to0.0484$), reproducing their component counts (79,
112, 149) exactly.

One limitation of that route is worth stating, because it is a modelling
assumption rather than an approximation that improves with effort. The
corresponding-states construction takes the valence to be PROPORTIONAL TO
SURFACE AREA, so charge polydispersity is fully correlated with size
polydispersity: one distribution governs both. A sample whose surface
chemistry does not scale with area --- differing grafting densities across a
size fraction, say --- cannot be represented.

\citet{GinozaYasutomi1998} treat size and interaction polydispersity as
\emph{independent}, on the same factorisable-coefficient mean-spherical
solution, and would be the route to take if that case ever arose. Both
descend from \citet{Ginoza1986}, who showed that factorisable coefficients
$K_{ij} = K d_i d_j$ collapse the multicomponent Blum--H\o ye equations to
rational functions of a single scalar $\Gamma$ --- which is what makes any of
this affordable for a size distribution, since each component then enters
only through $(\sigma_i, \rho_i, d_i)$ and the cost per evaluation is linear
in the number of classes rather than quadratic.

A partial implementation of the Ginoza route exists in
\code{ginoza\_yukawa\_msa.py} and is deliberately shelved: it solves for
$\Gamma$ and stops short of $S(q)$. Finishing it would reproduce, from
equations transcribed out of a scanned 1986 paper, a structure factor the
Gazzillo route already computes and which is validated above at machine
precision. The two closing conditions Ginoza gives, his Eqs.~(38a) and
(38b), are both implemented there and agree to $10^{-12}$, so the
transcription is at least self-consistent; that is recorded in case the
decoupled-charge case ever makes the module worth completing.

\paragraph{Polydisperse sticky hard spheres.}
\citet{GazzilloGiacometti2000}, a mean-spherical solution in the Baxter sticky
limit with factorisable coupling parameters $K_{ij}=Y_iY_j$, which renders
the Baxter factor matrix three-dyadic and collapses the $p\times p$ inversion
to a $4\times4$ determinant for any number of components. Validation:
the $\gamma_0=0$ limits against analytic Percus--Yevick ($4\times10^{-14}$)
and \code{mixscatter} ($1\times10^{-13}$); the monodisperse sticky
compressibility against their closed form ($3\times10^{-8}$); the critical
point $\eta_c=(\sqrt3-1)/2$, $\gamma_{0c}^2=(\sqrt3+2)/6$ located as the
divergence of $S(0)$; and the generalised Boyle temperatures of their
Eq.~(53) to the printed digits ($2.787$, $1.677$, $2.986$, $2.897$ against
$2.79$, $1.68$, $2.99$, $2.90$).

Note that the stickiness parameter of this model is \emph{not} Baxter's
$\tau$ \citep{Baxter1968}: Baxter's Percus--Yevick solution for mixtures requires $p(p+1)/2$
coupled quadratic equations and is not in general dyadic, which is why the
mean-spherical model exists. The two agree only in the hard-sphere limit.

Three entries in Table~\ref{tab:software} deserve comment.

The \textbf{RMSA} comparison was exact at all nine state points tested
($\eta = 0.05$--$0.25$, screening lengths $2$--$10$\,nm, contact potentials
$20$--$100\,k_BT$), including four at which the Hayter--Penfold rescaling
\citep{HayterPenfold1981}
procedure engaged, so the agreement is not merely that of the underlying MSA
branch. Two independently written implementations of a numerically delicate
algorithm reproduce each other to all printed digits.

The \textbf{two-Yukawa} comparison validates the hard-core double-Yukawa
potential, including cases of mixed sign (attractive $K_1$, repulsive $K_2$),
which is the regime in which a sign convention is most easily mistaken. The
reference implementation states the identical convention,
$V(r) = -k_BT[K_1 e^{-Z_1(r-1)} + K_2 e^{-Z_2(r-1)}]/r$ for $r>\sigma$.

The \textbf{numerical} comparison is the only external check available for the
hypernetted-chain closure, which has no analytic solution. Both solvers are
discretised and differ in transform, mixing scheme and hard-core
representation, so the $0.3\,\%$ agreement is the expected order.

\subsection{Two defects found by external comparison}

\paragraph{Identical hard cores in mixtures.}
Comparison with \code{mixscatter} revealed a discrepancy of $5.2\,\%$ that,
diagnostically, did \emph{not} decrease under grid refinement and was
\emph{linear in the size spread}, vanishing exactly in the monodisperse limit.
The cause was that the hard-sphere potential routine positioned its core by
array index rather than against the radial grid, so it did not see the
rescaled grid on which the polydisperse construction of
equation~(\ref{eq:reducedtail}) relies. Every pair therefore received its core
at the same radius: the mixture was being solved as though all particles were
identical. After correction the agreement is $0.09\,\%$ and converges. It was
the only one of twenty-one potentials written in that form.

\subsection{Two numerical traps in fitting}

Both of the following are properties of the fitting problem rather than of
this implementation, and both produce a converged fit with a plausible
$\chi^2$ and silently wrong parameters. We report them because neither is
evident from the literature and both cost us considerable time.

\paragraph{A rank-revealing linear solve can discard a parameter entirely.}
It is natural to split the fit: the scale, the flat background and any
power-law amplitude enter the intensity linearly and can be obtained exactly
by weighted least squares at every iteration, leaving only the shape
parameters to the nonlinear optimiser. The saving is real. The hazard is
that the design columns need not be comparable in magnitude. A model built
from scattering length densities in $\mathrm{cm}^{-2}$ carries the contrast
squared and so has magnitude $\sim10^{20}$, while the constant background
column is unity --- a condition number near $10^{20}$.
\code{numpy.linalg.lstsq}, like any rank-revealing solver, discards singular
values below $\max(M,N)\epsilon$, so the background and power-law
coefficients are returned as \emph{exactly} zero and the solve reports
success. The symptom is a fit that a hand-chosen background beats by tens of
percent in $\chi^2$, which is indistinguishable from a broken optimiser.
Equilibrating the columns to unit norm before the solve and rescaling
afterwards recovers all coefficients to five figures and changes nothing
where the problem was already well conditioned \citep{GolubVanLoan2013}.

The same scaling defeats the nonlinear optimiser: a parameter of order
$10^{10}$ satisfies \code{least\_squares}'s gradient test after a single
evaluation, so it never moves and is reported with a meaningless
uncertainty. Since the absolute scattering length densities are in any case
perfectly correlated with the scale, we fit their RATIO, which is of order
unity and whose sign is the quantity the data actually determine.

\paragraph{Gaussian resolution smearing fails where it is needed most.} A
small-angle instrument broadens $\mathbf{Q}$ in two dimensions in the
detector plane, and the data are then radially averaged; the distribution of
the magnitude about a true $Q_0$ is Rician, Eq.~\eqref{eq:rician} below, not
Gaussian \citep{Pedersen1990,BarkerPedersen1995}. A Gaussian in $Q$ is its
$Q_0/\sigma \gg 1$ limit and is entirely adequate wherever
$\mathrm{d}Q/Q$ is small.

\begin{equation}
  P(Q\,|\,Q_0) \;=\; \frac{Q}{\sigma^{2}}
  \exp\!\left(-\frac{Q^{2}+Q_0^{2}}{2\sigma^{2}}\right)
  I_0\!\left(\frac{QQ_0}{\sigma^{2}}\right)
  \label{eq:rician}
\end{equation}

It is not always small. For a spherical-shell dataset used to develop this
work, $\mathrm{d}Q/Q$ reaches $0.55$ at the lowest measured point, where a
Gaussian kernel places $44\,\%$ of that point's weight below $Q_{\min}$ and
some of it at \emph{negative} $Q$. The leading factor of $Q$ in
\eqref{eq:rician} suppresses the origin and excludes negative $Q$ by
construction, reducing that to $34\,\%$. Both forms reproduce the unsmeared
curve to $2\times10^{-16}$ as $\mathrm{d}Q \to 0$, which is the check that
catches a mis-normalised kernel. Evaluate \eqref{eq:rician} through the
exponentially scaled $\tilde{I}_0(x) = e^{-|x|}I_0(x)$: the argument reaches
several hundred at these widths, where $I_0$ overflows a double long before
the prefactor cancels it.

A related trap follows. A background term $A\,Q^{-4+d}$ diverges at the
origin while the smearing kernel reaches below the measured range, so the
smeared value at the lowest point is governed by whatever guard prevents
division by zero --- in our case $6.6\times10^{5}$ times the unsmeared value,
and varying by five orders of magnitude as that guard was moved. The power
law describes the measured range and says nothing about $Q \to 0$, where any
real curve turns over; holding it constant below the lowest measured $Q$
removes the arbitrariness.

\begin{quote}\small\itshape
Draft note on citations. Now placed: \code{OrnsteinZernike1914} and
\code{PercusYevick1958} (the equation and closure, \S2),
\code{HansenMcDonald2013}, \code{Anderson1965}, \code{WalkerNi2011},
\code{HayterPenfold1981}, \code{Liu2005}, \code{DiazMaier2024}
(\code{mixscatter}), \code{BarkerPedersen1995} and
\code{GolubVanLoan2013}.

Still in \code{ozgui.bib} and uncited: \code{DuhHenderson1996},
\code{Gazzillo1999}, \code{Gazzillo2006}, \code{KalyuzhnyiCummings2006},
\code{Carsughi2000}. Each has a natural home --- the first in the discussion
of mixing rules, the Gazzillo papers where stickiness and its size
dependence are treated --- but placing a citation where it is not actually
used is worse than leaving it out, so they are left for whoever writes those
sections.

Still absent and worth adding: software citations for NumPy, SciPy, SUNDIALS
and NLopt, all of which the work depends on.

One substantive point from the verification. \citet{BarkerPedersen1995}
conclude that a Gaussian approximation to the resolution is PREFERRED in most
experiments for ease of calculation, and is inadequate only in particular
cases --- of which they name Porod scattering. The argument above should be
framed as identifying one of those known exceptions rather than as a general
correction, which is both more defensible and more useful.
\end{quote}

\subsection{What these figures measure}

A note on what these figures measure. Every agreement quoted in this section
was obtained with the type-1 discrete sine transform, whose grid points fall
\emph{on} the hard core, and which is therefore first order in $\Delta r$. A
type-4 transform places the grid at $(n+\tfrac12)\Delta r$ so that the core
lies between points --- the midpoint rule --- and is second order. For one
component it is both more accurate and cheaper, and both advantages grow with
resolution:

\begin{center}\small
\begin{tabular}{rrrr}
\toprule
points per $\sigma$ & type-4 time & accuracy gain\\
\midrule
100 & $0.73\times$ & $139\times$\\
400 & $0.65\times$ & $547\times$\\
800 & $0.49\times$ & $1092\times$\\
\bottomrule
\end{tabular}
\end{center}

The times are relative to type 1 at the same points per diameter; the
accuracy gain is the ratio of the two errors. The growth of that ratio with
resolution is the second-order against first-order signature. The figures
quoted above would improve correspondingly.

The two transforms require different grid \emph{lengths}, which is easy to
overlook and costly to get wrong: a DST-I is computed through an FFT of
length $2(N+1)$ and so wants $N = 2^{k}-1$, whereas a DST-IV is a length-$N$
transform and wants $N = 2^{k}$. Sizing both for the former leaves the latter
on a composite length ($16383 = 3\times43\times127$) and costs it a factor of
three in the transform, enough to make the second-order method appear slower
than the first-order one. Accuracy is unaffected by the choice.

We nevertheless retain type 1 as the default, because the type-4 condition
cannot be met for a mixture. It requires every pair core $\sigma_{ij}$ to
fall between grid points, and with size classes drawn from a Gaussian
quadrature the $\sigma_i$ are irrational, so the $p(p+1)/2$ conditions cannot
be satisfied simultaneously by any choice of resolution. Type 4 is available
as an option and is correct for one component, where the single core
coincides with the reduced unit by construction.


The defect had been invisible to every existing test because the resulting
error vanishes in precisely the monodisperse limit where those tests looked.

\paragraph{Under-resolved narrow wells.}
\label{sec:welltrap}
Comparison with an analytic adhesive-sphere solution showed errors that are
governed not by the number of grid points per particle diameter but by the
number \emph{across the attractive well}:

\begin{center}\small
\begin{tabular}{cc}
\toprule
points across well & max.\ relative error\\
\midrule
2 & 0.916\\
4 & 0.324\\
8 & 0.142\\
16 & 0.068\\
32 & 0.035\\
\bottomrule
\end{tabular}
\end{center}

This is ordinary first-order convergence in $\Delta r/\delta$ and not a
defect, but the default discretisation violates it badly: at 100 points per
diameter a well of width $\delta = 0.02\sigma$ spans two grid points and the
structure factor is in error by ninety per cent, without any warning. We
recommend $\;\gtrsim 30\,\sigma/\delta$ points per diameter for any potential
with a narrow feature, and note that this requirement applies equally to
square-well and piecewise-constant potentials.

\subsection{Internal consistency}

The following are not independent checks but constrain the implementation:
reduction of the polydisperse construction to a single size class reproduces
the original one-component routines bit-identically; the extended
Rogers--Young closure reduces to Rogers--Young to $2.5\times10^{-16}$ when its
additional parameter vanishes; the hard-core double-Yukawa potential agrees
with the three-Yukawa form to $\sim10^{-16}$ under the stated
reparameterisation; and the exact intensity reproduces brute-force
integration of $\langle|F|^2\rangle$ in the dilute limit
(ripple $0.041$ against $0.038$).

% ---------------------------------------------------------------------
% References to add to the bibliography
% ---------------------------------------------------------------------
% Gazzillo, D., Giacometti, A. & Carsughi, F. (1997). J. Chem. Phys. 107,
%   10141. [polydisperse charged hard spheres, MSA]
% Gazzillo, D., Giacometti, A. & Carsughi, F. (1999). arXiv:cond-mat/9909153.
%   [corrected Appendix A; scaling vs decoupling approximation]
% Gazzillo, D. & Giacometti, A. (2000). arXiv:cond-mat/0011330.
%   [polydisperse sticky hard spheres, MSA, factorisable coupling]
% Carsughi, F., Giacometti, A. & Gazzillo, D. (2000). arXiv:cond-mat/0011242.
%   [the FLAC program]
% Tutschka, C. & Kahl, G. (1998). J. Chem. Phys. 108, 9498.
% Blum, L. & Hoye, J. S. (1978). J. Stat. Phys. 19, 317.
%   [multicomponent multi-Yukawa MSA; the finite-range generalisation]
% Hoye, J. S. & Stell, G. [SCOZA]; Reiner, A. & Hoye, J. S. (2008).
%   arXiv:0712.3739. [SCOZA near the critical point]
% Bergmann, M. J. et al. mixscatter (Python package).
%   [add version/DOI; used for the analytic Vrij mixture PY]
% Biehl, R. (2019). jscatter. PLoS ONE 14, e0218789.
%   [Two_Yukawa after Liu et al. (2005); RMSA after Hayter & Penfold (1981)]
% Hayter, J. B. & Penfold, J. (1981). Mol. Phys. 42, 109-118.
% Liu, Y., Chen, W.-R. & Chen, S.-H. (2005). J. Chem. Phys. 122, 044507.
%   [two-Yukawa analytic solution]
% Pihlajamaa, I. & Janssen, L. M. C. (2024). Phys. Rev. E 110, 044608.
%   [closure comparison; OrnsteinZernike.jl]
% Vrij, A. (1979). J. Chem. Phys. 71, 3267-3270.
%   [mixture PY, the reference for the mixture comparison]
% Wertheim, M. S. (1963). Phys. Rev. Lett. 10, 321.
% Thiele, E. (1963). J. Chem. Phys. 39, 474.
% Baxter, R. J. (1968). J. Chem. Phys. 49, 2770.
% Malijevsky, A. & Labik, S. (1987). Mol. Phys. 60, 663.
.
% Not standalone: no preamble, no \begin{document}.
% Every number here was measured against an INDEPENDENT implementation
% or an analytic result -- never against a previous version of this code.

\section{Validation}
\label{sec:validation}

The implementation was checked against analytic solutions, against published
tabulations, and against four independently written software packages. We
report the comparisons in full, including the two cases in which they
uncovered defects in our own code, because the pattern is itself informative:
every defect found late in this work was invisible to internal consistency
checks and was exposed only by an external reference.

\subsection{Analytic and literature references}

\begin{table}[htbp]
\centering
\caption{Comparison with analytic results and published tabulations. Where a
difference is quoted it is the deviation from the reference, not from a
previous version of this code.}
\label{tab:analytic}
\small
\begin{tabular}{llll}
\toprule
Reference & Quantity & Ours & Agreement\\
\midrule
\citet{Wertheim1963}, \citet{Thiele1963} & $c(r)$, PY hard spheres & numerical & first order in $\Delta r$\\
                               & $g(\sigma^+)$, $\eta=0.1$--$0.4$ & numerical & $0.1\,\%$\\
                               & $S(0)=(1-\eta)^4/(1+2\eta)^2$ & numerical & converges\\
Analytic PY, $\eta=0.3$        & $\chi^{-1}$ & 11.07 & 10.66 (extrapolation)\\
                               & $\beta P/\rho$ & 3.890 & 3.816 (grid)\\
                               & PY inconsistency & $+1.33$ & $+1.25$\\
D'Aguanno \& Klein (1991)      & $1/S(0)$, excess pressure & --- & $0.01\,\%$, $0.02\,\%$\\
                               & excess energy, $g_{\max}$ & --- & $0.27\,\%$, $0.54\,\%$\\
\citet{Vrij1979} mixture PY         & $S^{\mathrm{AL}}_{ij}$, 3 classes & numerical & $0.09\,\%$\\
\bottomrule
\end{tabular}
\end{table}

\subsection{Comparison with independent software}

\begin{table}[htbp]
\centering
\caption{Comparison with independently written packages. ``exact'' denotes
agreement to all printed digits; the remaining entries are limited by
discretisation and converge as $\Delta r$ is reduced.}
\label{tab:software}
\small
\begin{tabular}{lllc}
\toprule
Package & Model & Method & Agreement\\
\midrule
\code{mixscatter} \citep{DiazMaier2024} & mixture PY \citep{Vrij1979} & analytic & $0.09\,\%$\\
\code{jscatter} & one-component PY & analytic & $0.02$--$0.12\,\%$\\
\code{sasmodels}, \code{jscatter} & RMSA \citep{HayterPenfold1981} & analytic & $10^{-12}$ (6/6)\\
\code{sasmodels} & two-Yukawa MSA \citep{Liu2005} & analytic & $0.21\,\%$\\
\code{jscatter} & adhesive hard spheres & analytic & see \S\ref{sec:welltrap}\\
\code{jscatter} & OZ solver, PY and HNC & numerical & $0.2$--$0.4\,\%$\\
\code{OrnsteinZernike.jl} & 4 bridge functions & analytic & $\sim10^{-16}$\\
\bottomrule
\end{tabular}
\end{table}

\paragraph{An absolute check: the compressibility.} Every comparison above
sets one calculation against another. One quantity admits an exact answer
instead. For monodisperse hard spheres the Percus--Yevick compressibility is
\begin{equation}
  S(0) = \frac{(1-\varphi)^{4}}{(1+2\varphi)^{2}},
  \label{eq:pyS0}
\end{equation}
so extrapolating the computed partials to the origin and comparing measures
absolute accuracy rather than agreement.

The result is that the transform type matters far more than its place in a
cost table suggests. At $\varphi = 0.40$ and 100 radial points per diameter,
the first-order DST-I gives $S(0)$ low by $5.27\,\%$; the second-order DST-IV
is low by $0.036\,\%$ on the same grid, and by $0.004\,\%$ at four times the
resolution. Type 1 refines as $3.96$ per fourfold step, textbook first order,
so it reaches $1.33\,\%$ where type 4 began 150 times better.

The cause is where the discontinuity lands. A DST-I places nodes at
$r = (n+1)\Delta r$, so the hard core at $r = \sigma$ falls ON a node and is
carried by a point that is neither inside nor outside; the excluded volume is
then wrong at first order in $\Delta r$. A DST-IV places them at
$(n+\frac{1}{2})\Delta r$, the core falls between points, and the step is
resolved to second order.

Two things make this worth reporting rather than absorbing into a
recommendation. First, the error is independent of the real-space range:
extending $r_{\max}$ from $41\sigma$ to $655\sigma$ left $S(0)$ unchanged in
the sixth decimal. A convergence study that refines resolution sees orderly
first-order behaviour and concludes the method works --- which it does, slowly,
towards an answer whose leading error is set by grid alignment. Second,
$S(0)$ is the isothermal compressibility, so this is a systematic bias in
exactly the quantity the volume fraction controls, and a fit will move
$\varphi$ to absorb it.

It also accounts for a discrepancy that had resisted explanation: a
multicomponent solve converged to a fixed point at $7\times10^{-12}$ yet
differing from two independent analytic references by $\sim10^{-3}$, the
difference immune to refinement and growing with $\varphi$. That is the same
effect seen from the other side.

The practical consequence is a restriction rather than a fix. Type 4 needs
every pair core between grid points, which one size class can be given and a
mixture cannot: the pair diameters $\sigma_{ij} = (\sigma_i+\sigma_j)/2$ come
from a moment-matched quadrature and are mutually incommensurate. The
polydisperse route is therefore first order, and until a mixture-capable
second-order treatment exists a dense polydisperse fit should use the finest
affordable resolution, with its fitted volume fraction read accordingly.

\paragraph{Reproducibility of the figures and of the table.} Three gaps are
recorded here rather than left to be discovered.

The remaining \code{jscatter} rows above --- one-component PY, adhesive
hard spheres and the numerical OZ comparison --- were measured where
\code{jscatter} was installed. Its dependency chain does not build under
MinGW, so they cannot be re-run on the machine this work is written on. Two
rows have since been brought back within reach, and by the same move:
replacing a reference that has to be obtained with one that can be
installed.

\paragraph{Two reference solvers replaced by \code{sasmodels}.} The RMSA
and two-Yukawa rows are now measured against SasView's \code{sasmodels},
which is 3-clause BSD and installs with \code{pip} alone. Both previously
depended on something a reader could not simply obtain --- \code{jscatter},
which will not build here, and a 97\,kB vendored copy of the two-Yukawa
solver used under a citation-ware grant the authors gave in writing. A
licence governs an implementation rather than a method, so \citet{Liu2005}
is cited exactly as before; what changes is that anyone can now repeat the
comparison.

Agreement is to $2\times10^{-8}$, $7\times10^{-8}$ and $3\times10^{-8}$ on
the three two-Yukawa state points where both solve cleanly. Of the
remaining two, \code{sasmodels} returns a physical $S(q)$ at
$\varphi = 0.10$ where the vendored solver reports no root found --- a
root-finder limitation rather than a statement about the physics, since the
same parameters at HIGHER volume fraction solve in both. The last,
$Z_1 = 6$ against $Z_2 = 4$, agrees only to $2\times10^{-3}$, with the
discrepancy peaking at the structure-factor maximum; the two inverse ranges
are close there, which is where the quartic's roots crowd together and two
implementations may select marginally differently. Both results are
physical and both are defensible.

\paragraph{The sign of $K$, and where the inconsistency actually was.}
\citet{Liu2005} write the tail as $-K\exp[-Z(r-1)]/r$, so positive $K$ is
ATTRACTIVE. This package's solver follows that convention, and so does
\code{sasmodels}. The vendored reference file did not: reproducing its
output required flipping both signs, after which the two agree to
$10^{-8}$, where every other candidate (length unit, pair ordering, $Z$
scaling) stayed above 20 per cent at every state point.

So the inconsistency was in the comparison apparatus rather than in the
solver being compared --- which is the more flattering of the two
possibilities, and was not the one assumed while the discrepancy was being
chased. Retiring the vendored file removes it entirely, and is a further
reason to prefer \code{sasmodels}: one fewer convention in the tree.

\code{fig3\_oscillation.pdf} has no generator in the repository. The figure
exists and nothing can reproduce it, which means a referee asking for any
change to it would find there is nothing to change.
\code{fig\_survey.pdf} names \code{tools/section53\_survey.py} in its own
caption and is the model the others should follow;
\code{fig1\_discretisation.pdf} now has
\code{tools/fig1\_discretisation.py}.

Regenerating it also caught a discrepancy in the figure itself.
\code{polydisperse\_nodes} offers two discretisations --- \code{sizeClasses},
the moment-matched generalised Gauss--Laguerre rule the caption describes,
and \code{quantileClasses}, which is quantile-based --- and only the first
reproduces the moments exactly. At $p = 3$ the Gauss--Laguerre rule returns
$\langle\sigma\rangle = 1.000000$ and $\langle\sigma^{2}\rangle = 1.062500$,
exactly $1 + s^{2}$; the quantile rule returns $0.9794$ and $0.9596$. Both
are usable discretisations, but only the first supports the claim printed
beside the figure. The moments are now shown on each panel so that the claim
can be checked rather than taken on trust.

\paragraph{The table's figures depend on the transform, and the table does
not say so.} The one-component Percus--Yevick row is quoted as
$0.02$--$0.12\,\%$ over $0.2 \le \varphi \le 0.4$, and that is a DST-IV
result. With the first-order DST-I at the same 100 radial points per
diameter the same comparison gives $0.70$, $1.37$ and $2.33\,\%$ at
$\varphi = 0.2$, $0.3$ and $0.4$ --- an order of magnitude worse, rising with
density, and reducible only fourfold per fourfold refinement
($0.17$, $0.33$, $0.51\,\%$ at 400 points).

This is the grid-alignment effect described above --- the hard core falling ON
a DST-I node rather than between DST-IV nodes --- seen in a second
quantity, and it carries the same caveat: type 4 is available for ONE size
class and not for a mixture, so the polydisperse route --- which is what this
paper is about --- carries the first-order figures rather than the quoted
ones. A reader comparing packages should know which transform produced a
number before comparing it with their own.

\paragraph{RMSA, and a units error that imitated a physics disagreement.}
\label{sec:rmsarescale}
This row agrees with \emph{two} independent implementations to $10^{-12}$
over $\gamma = 3.7$ to $816$ and $\kappa\sigma = 1.5$ to $3.8$, including the
strongly coupled points at which the Hayter--Penfold rescaling engages.
Reaching that took an afternoon, and the detour is worth recording because
the symptom was convincingly misleading.

The earlier claim was ``exact (9/9)'' against \code{jscatter}, which does not
build under MinGW: unrepeatable on the machine this is developed on.
\code{sasmodels}' \code{hayter\_msa} is a better reference --- a different
code lineage, installable with \code{pip} and nothing else --- but comparing
against it gave disagreements of 6 to 32 per cent that GREW WITH COUPLING.
That is the signature of a physics difference rather than a units slip, and
it was pursued as one: root selection between the two solutions the closure
admits, the rescaling threshold, radius-versus-diameter in $\gamma$. Each
was plausible, each was testable, each was wrong. No single rescaling of any
input reconciled the two, and the best-fitting factor itself drifted with
coupling.

The answer is in Eq.~(3a) of \citet{HayterPenfold1981}. The reduced pair
potential is written
\begin{equation}
\beta U(x) = \gamma\,\frac{e^{-kx}}{x},
\qquad x = r/\sigma,\quad k = \kappa\sigma,
\label{eq:hpgamma}
\end{equation}
so the value AT CONTACT is $\gamma e^{-k}$, and that --- not $\gamma$ --- is
what the paper equates to $\beta z^2/[\pi\epsilon_0\epsilon\sigma(2+k)^2]$.
\code{sasmodels} computes the contact potential; \code{jscatter} and the
wrapper here take the amplitude $\gamma$. The two differ by $e^{k}$, and
because $k$ varies with the screening the discrepancy grows with coupling
--- which is exactly why it read as physics.

Two things are worth taking from this. A discrepancy that scales with a
physical parameter is NOT evidence that the difference is physical: a
conversion factor depending on that parameter behaves identically. And the
arbiter was the original paper rather than either code's documentation ---
both describe $\gamma$ as ``the contact potential'', and both use the phrase
loosely for a quantity their own sources define differently.

With the factor applied, \code{tools/rmsa\_vs\_sasmodels.py} reproduces
\code{sasmodels} at all six state points, and the same inputs reproduce
\code{jscatter} to the same precision. The claim in the table is therefore
checkable by any reader with \code{pip install sasmodels}, which the one it
replaces was not.

\paragraph{A closure whose domain is unmapped.} The Carbajal--Tinoco bridge
function \citep{CarbajalTinoco2022} is implicit, $b = T(\gamma + b)$, and
has no solution at all below a
threshold: substituting $w = \gamma + b$ turns it into the explicit map
$\gamma(w) = w - T(w)$, whose range is bounded below. Numerically
$\gamma_{\min} = -0.509$ at $\lambda = 0$, rising towards zero as $\lambda$
grows, with two solutions above it and none below. Neither the source paper
nor the comparison of \citet{PihlajamaaJanssen2024} locates that boundary;
the latter simply excludes closures from its figures ``due to convergence
issues at this state point''.

\subsection{Analytic references implemented in this work}
\label{sec:analyticrefs}

Two of the three classical solvable polydisperse models had no reference
implementation available to us, so we implemented them. Together with the
Vrij hard-sphere solution supplied by \code{mixscatter}, they give exact
analytic coverage of the polydisperse hard-sphere, charged-hard-sphere and
sticky-hard-sphere fluids, against which the numerical solver can be checked
at any state point.

\paragraph{Polydisperse charged hard spheres.}
Appendix A of Gazzillo, Giacometti \& Carsughi (1999), the corrected form of
their 1997 paper --- the correction supplies a factor $\pi$ omitted from
$P_z$ and repairs an error in the determinant. This is the same physics
implemented by the FLAC program of Carsughi \emph{et al.}, which cites the
uncorrected 1997 paper. Our implementation reproduces the neutral
monodisperse limit against analytic Percus--Yevick to $3\times10^{-14}$, the
neutral polydisperse limit against \code{mixscatter} to $1\times10^{-13}$,
and, for charged systems, all three effects the authors report for increasing
polydispersity: the first peak is lowered ($1.7350\to1.3546$), shifted to
smaller $Q$ ($5.841\to5.030$ in units of $Q\langle\sigma\rangle$) and
$S(0)$ increased ($0.0187\to0.0484$), reproducing their component counts (79,
112, 149) exactly.

One limitation of that route is worth stating, because it is a modelling
assumption rather than an approximation that improves with effort. The
corresponding-states construction takes the valence to be PROPORTIONAL TO
SURFACE AREA, so charge polydispersity is fully correlated with size
polydispersity: one distribution governs both. A sample whose surface
chemistry does not scale with area --- differing grafting densities across a
size fraction, say --- cannot be represented.

\citet{GinozaYasutomi1998} treat size and interaction polydispersity as
\emph{independent}, on the same factorisable-coefficient mean-spherical
solution, and would be the route to take if that case ever arose. Both
descend from \citet{Ginoza1986}, who showed that factorisable coefficients
$K_{ij} = K d_i d_j$ collapse the multicomponent Blum--H\o ye equations to
rational functions of a single scalar $\Gamma$ --- which is what makes any of
this affordable for a size distribution, since each component then enters
only through $(\sigma_i, \rho_i, d_i)$ and the cost per evaluation is linear
in the number of classes rather than quadratic.

A partial implementation of the Ginoza route exists in
\code{ginoza\_yukawa\_msa.py} and is deliberately shelved: it solves for
$\Gamma$ and stops short of $S(q)$. Finishing it would reproduce, from
equations transcribed out of a scanned 1986 paper, a structure factor the
Gazzillo route already computes and which is validated above at machine
precision. The two closing conditions Ginoza gives, his Eqs.~(38a) and
(38b), are both implemented there and agree to $10^{-12}$, so the
transcription is at least self-consistent; that is recorded in case the
decoupled-charge case ever makes the module worth completing.

\paragraph{Polydisperse sticky hard spheres.}
\citet{GazzilloGiacometti2000}, a mean-spherical solution in the Baxter sticky
limit with factorisable coupling parameters $K_{ij}=Y_iY_j$, which renders
the Baxter factor matrix three-dyadic and collapses the $p\times p$ inversion
to a $4\times4$ determinant for any number of components. Validation:
the $\gamma_0=0$ limits against analytic Percus--Yevick ($4\times10^{-14}$)
and \code{mixscatter} ($1\times10^{-13}$); the monodisperse sticky
compressibility against their closed form ($3\times10^{-8}$); the critical
point $\eta_c=(\sqrt3-1)/2$, $\gamma_{0c}^2=(\sqrt3+2)/6$ located as the
divergence of $S(0)$; and the generalised Boyle temperatures of their
Eq.~(53) to the printed digits ($2.787$, $1.677$, $2.986$, $2.897$ against
$2.79$, $1.68$, $2.99$, $2.90$).

Note that the stickiness parameter of this model is \emph{not} Baxter's
$\tau$ \citep{Baxter1968}: Baxter's Percus--Yevick solution for mixtures requires $p(p+1)/2$
coupled quadratic equations and is not in general dyadic, which is why the
mean-spherical model exists. The two agree only in the hard-sphere limit.

Three entries in Table~\ref{tab:software} deserve comment.

The \textbf{RMSA} comparison was exact at all nine state points tested
($\eta = 0.05$--$0.25$, screening lengths $2$--$10$\,nm, contact potentials
$20$--$100\,k_BT$), including four at which the Hayter--Penfold rescaling
\citep{HayterPenfold1981}
procedure engaged, so the agreement is not merely that of the underlying MSA
branch. Two independently written implementations of a numerically delicate
algorithm reproduce each other to all printed digits.

The \textbf{two-Yukawa} comparison validates the hard-core double-Yukawa
potential, including cases of mixed sign (attractive $K_1$, repulsive $K_2$),
which is the regime in which a sign convention is most easily mistaken. The
reference implementation states the identical convention,
$V(r) = -k_BT[K_1 e^{-Z_1(r-1)} + K_2 e^{-Z_2(r-1)}]/r$ for $r>\sigma$.

The \textbf{numerical} comparison is the only external check available for the
hypernetted-chain closure, which has no analytic solution. Both solvers are
discretised and differ in transform, mixing scheme and hard-core
representation, so the $0.3\,\%$ agreement is the expected order.

\subsection{Two defects found by external comparison}

\paragraph{Identical hard cores in mixtures.}
Comparison with \code{mixscatter} revealed a discrepancy of $5.2\,\%$ that,
diagnostically, did \emph{not} decrease under grid refinement and was
\emph{linear in the size spread}, vanishing exactly in the monodisperse limit.
The cause was that the hard-sphere potential routine positioned its core by
array index rather than against the radial grid, so it did not see the
rescaled grid on which the polydisperse construction of
equation~(\ref{eq:reducedtail}) relies. Every pair therefore received its core
at the same radius: the mixture was being solved as though all particles were
identical. After correction the agreement is $0.09\,\%$ and converges. It was
the only one of twenty-one potentials written in that form.

\subsection{Two numerical traps in fitting}

Both of the following are properties of the fitting problem rather than of
this implementation, and both produce a converged fit with a plausible
$\chi^2$ and silently wrong parameters. We report them because neither is
evident from the literature and both cost us considerable time.

\paragraph{A rank-revealing linear solve can discard a parameter entirely.}
It is natural to split the fit: the scale, the flat background and any
power-law amplitude enter the intensity linearly and can be obtained exactly
by weighted least squares at every iteration, leaving only the shape
parameters to the nonlinear optimiser. The saving is real. The hazard is
that the design columns need not be comparable in magnitude. A model built
from scattering length densities in $\mathrm{cm}^{-2}$ carries the contrast
squared and so has magnitude $\sim10^{20}$, while the constant background
column is unity --- a condition number near $10^{20}$.
\code{numpy.linalg.lstsq}, like any rank-revealing solver, discards singular
values below $\max(M,N)\epsilon$, so the background and power-law
coefficients are returned as \emph{exactly} zero and the solve reports
success. The symptom is a fit that a hand-chosen background beats by tens of
percent in $\chi^2$, which is indistinguishable from a broken optimiser.
Equilibrating the columns to unit norm before the solve and rescaling
afterwards recovers all coefficients to five figures and changes nothing
where the problem was already well conditioned \citep{GolubVanLoan2013}.

The same scaling defeats the nonlinear optimiser: a parameter of order
$10^{10}$ satisfies \code{least\_squares}'s gradient test after a single
evaluation, so it never moves and is reported with a meaningless
uncertainty. Since the absolute scattering length densities are in any case
perfectly correlated with the scale, we fit their RATIO, which is of order
unity and whose sign is the quantity the data actually determine.

\paragraph{Gaussian resolution smearing fails where it is needed most.} A
small-angle instrument broadens $\mathbf{Q}$ in two dimensions in the
detector plane, and the data are then radially averaged; the distribution of
the magnitude about a true $Q_0$ is Rician, Eq.~\eqref{eq:rician} below, not
Gaussian \citep{Pedersen1990,BarkerPedersen1995}. A Gaussian in $Q$ is its
$Q_0/\sigma \gg 1$ limit and is entirely adequate wherever
$\mathrm{d}Q/Q$ is small.

\begin{equation}
  P(Q\,|\,Q_0) \;=\; \frac{Q}{\sigma^{2}}
  \exp\!\left(-\frac{Q^{2}+Q_0^{2}}{2\sigma^{2}}\right)
  I_0\!\left(\frac{QQ_0}{\sigma^{2}}\right)
  \label{eq:rician}
\end{equation}

It is not always small. For a spherical-shell dataset used to develop this
work, $\mathrm{d}Q/Q$ reaches $0.55$ at the lowest measured point, where a
Gaussian kernel places $44\,\%$ of that point's weight below $Q_{\min}$ and
some of it at \emph{negative} $Q$. The leading factor of $Q$ in
\eqref{eq:rician} suppresses the origin and excludes negative $Q$ by
construction, reducing that to $34\,\%$. Both forms reproduce the unsmeared
curve to $2\times10^{-16}$ as $\mathrm{d}Q \to 0$, which is the check that
catches a mis-normalised kernel. Evaluate \eqref{eq:rician} through the
exponentially scaled $\tilde{I}_0(x) = e^{-|x|}I_0(x)$: the argument reaches
several hundred at these widths, where $I_0$ overflows a double long before
the prefactor cancels it.

A related trap follows. A background term $A\,Q^{-4+d}$ diverges at the
origin while the smearing kernel reaches below the measured range, so the
smeared value at the lowest point is governed by whatever guard prevents
division by zero --- in our case $6.6\times10^{5}$ times the unsmeared value,
and varying by five orders of magnitude as that guard was moved. The power
law describes the measured range and says nothing about $Q \to 0$, where any
real curve turns over; holding it constant below the lowest measured $Q$
removes the arbitrariness.

\begin{quote}\small\itshape
Draft note on citations. Now placed: \code{OrnsteinZernike1914} and
\code{PercusYevick1958} (the equation and closure, \S2),
\code{HansenMcDonald2013}, \code{Anderson1965}, \code{WalkerNi2011},
\code{HayterPenfold1981}, \code{Liu2005}, \code{DiazMaier2024}
(\code{mixscatter}), \code{BarkerPedersen1995} and
\code{GolubVanLoan2013}.

Still in \code{ozgui.bib} and uncited: \code{DuhHenderson1996},
\code{Gazzillo1999}, \code{Gazzillo2006}, \code{KalyuzhnyiCummings2006},
\code{Carsughi2000}. Each has a natural home --- the first in the discussion
of mixing rules, the Gazzillo papers where stickiness and its size
dependence are treated --- but placing a citation where it is not actually
used is worse than leaving it out, so they are left for whoever writes those
sections.

Still absent and worth adding: software citations for NumPy, SciPy, SUNDIALS
and NLopt, all of which the work depends on.

One substantive point from the verification. \citet{BarkerPedersen1995}
conclude that a Gaussian approximation to the resolution is PREFERRED in most
experiments for ease of calculation, and is inadequate only in particular
cases --- of which they name Porod scattering. The argument above should be
framed as identifying one of those known exceptions rather than as a general
correction, which is both more defensible and more useful.
\end{quote}

\subsection{What these figures measure}

A note on what these figures measure. Every agreement quoted in this section
was obtained with the type-1 discrete sine transform, whose grid points fall
\emph{on} the hard core, and which is therefore first order in $\Delta r$. A
type-4 transform places the grid at $(n+\tfrac12)\Delta r$ so that the core
lies between points --- the midpoint rule --- and is second order. For one
component it is both more accurate and cheaper, and both advantages grow with
resolution:

\begin{center}\small
\begin{tabular}{rrrr}
\toprule
points per $\sigma$ & type-4 time & accuracy gain\\
\midrule
100 & $0.73\times$ & $139\times$\\
400 & $0.65\times$ & $547\times$\\
800 & $0.49\times$ & $1092\times$\\
\bottomrule
\end{tabular}
\end{center}

The times are relative to type 1 at the same points per diameter; the
accuracy gain is the ratio of the two errors. The growth of that ratio with
resolution is the second-order against first-order signature. The figures
quoted above would improve correspondingly.

The two transforms require different grid \emph{lengths}, which is easy to
overlook and costly to get wrong: a DST-I is computed through an FFT of
length $2(N+1)$ and so wants $N = 2^{k}-1$, whereas a DST-IV is a length-$N$
transform and wants $N = 2^{k}$. Sizing both for the former leaves the latter
on a composite length ($16383 = 3\times43\times127$) and costs it a factor of
three in the transform, enough to make the second-order method appear slower
than the first-order one. Accuracy is unaffected by the choice.

We nevertheless retain type 1 as the default, because the type-4 condition
cannot be met for a mixture. It requires every pair core $\sigma_{ij}$ to
fall between grid points, and with size classes drawn from a Gaussian
quadrature the $\sigma_i$ are irrational, so the $p(p+1)/2$ conditions cannot
be satisfied simultaneously by any choice of resolution. Type 4 is available
as an option and is correct for one component, where the single core
coincides with the reduced unit by construction.


The defect had been invisible to every existing test because the resulting
error vanishes in precisely the monodisperse limit where those tests looked.

\paragraph{Under-resolved narrow wells.}
\label{sec:welltrap}
Comparison with an analytic adhesive-sphere solution showed errors that are
governed not by the number of grid points per particle diameter but by the
number \emph{across the attractive well}:

\begin{center}\small
\begin{tabular}{cc}
\toprule
points across well & max.\ relative error\\
\midrule
2 & 0.916\\
4 & 0.324\\
8 & 0.142\\
16 & 0.068\\
32 & 0.035\\
\bottomrule
\end{tabular}
\end{center}

This is ordinary first-order convergence in $\Delta r/\delta$ and not a
defect, but the default discretisation violates it badly: at 100 points per
diameter a well of width $\delta = 0.02\sigma$ spans two grid points and the
structure factor is in error by ninety per cent, without any warning. We
recommend $\;\gtrsim 30\,\sigma/\delta$ points per diameter for any potential
with a narrow feature, and note that this requirement applies equally to
square-well and piecewise-constant potentials.

\subsection{Internal consistency}

The following are not independent checks but constrain the implementation:
reduction of the polydisperse construction to a single size class reproduces
the original one-component routines bit-identically; the extended
Rogers--Young closure reduces to Rogers--Young to $2.5\times10^{-16}$ when its
additional parameter vanishes; the hard-core double-Yukawa potential agrees
with the three-Yukawa form to $\sim10^{-16}$ under the stated
reparameterisation; and the exact intensity reproduces brute-force
integration of $\langle|F|^2\rangle$ in the dilute limit
(ripple $0.041$ against $0.038$).

% ---------------------------------------------------------------------
% References to add to the bibliography
% ---------------------------------------------------------------------
% Gazzillo, D., Giacometti, A. & Carsughi, F. (1997). J. Chem. Phys. 107,
%   10141. [polydisperse charged hard spheres, MSA]
% Gazzillo, D., Giacometti, A. & Carsughi, F. (1999). arXiv:cond-mat/9909153.
%   [corrected Appendix A; scaling vs decoupling approximation]
% Gazzillo, D. & Giacometti, A. (2000). arXiv:cond-mat/0011330.
%   [polydisperse sticky hard spheres, MSA, factorisable coupling]
% Carsughi, F., Giacometti, A. & Gazzillo, D. (2000). arXiv:cond-mat/0011242.
%   [the FLAC program]
% Tutschka, C. & Kahl, G. (1998). J. Chem. Phys. 108, 9498.
% Blum, L. & Hoye, J. S. (1978). J. Stat. Phys. 19, 317.
%   [multicomponent multi-Yukawa MSA; the finite-range generalisation]
% Hoye, J. S. & Stell, G. [SCOZA]; Reiner, A. & Hoye, J. S. (2008).
%   arXiv:0712.3739. [SCOZA near the critical point]
% Bergmann, M. J. et al. mixscatter (Python package).
%   [add version/DOI; used for the analytic Vrij mixture PY]
% Biehl, R. (2019). jscatter. PLoS ONE 14, e0218789.
%   [Two_Yukawa after Liu et al. (2005); RMSA after Hayter & Penfold (1981)]
% Hayter, J. B. & Penfold, J. (1981). Mol. Phys. 42, 109-118.
% Liu, Y., Chen, W.-R. & Chen, S.-H. (2005). J. Chem. Phys. 122, 044507.
%   [two-Yukawa analytic solution]
% Pihlajamaa, I. & Janssen, L. M. C. (2024). Phys. Rev. E 110, 044608.
%   [closure comparison; OrnsteinZernike.jl]
% Vrij, A. (1979). J. Chem. Phys. 71, 3267-3270.
%   [mixture PY, the reference for the mixture comparison]
% Wertheim, M. S. (1963). Phys. Rev. Lett. 10, 321.
% Thiele, E. (1963). J. Chem. Phys. 39, 474.
% Baxter, R. J. (1968). J. Chem. Phys. 49, 2770.
% Malijevsky, A. & Labik, S. (1987). Mol. Phys. 60, 663.
.
% Not standalone: no preamble, no \begin{document}.
% Every number here was measured against an INDEPENDENT implementation
% or an analytic result -- never against a previous version of this code.

\section{Validation}
\label{sec:validation}

The implementation was checked against analytic solutions, against published
tabulations, and against four independently written software packages. We
report the comparisons in full, including the two cases in which they
uncovered defects in our own code, because the pattern is itself informative:
every defect found late in this work was invisible to internal consistency
checks and was exposed only by an external reference.

\subsection{Analytic and literature references}

\begin{table}[htbp]
\centering
\caption{Comparison with analytic results and published tabulations. Where a
difference is quoted it is the deviation from the reference, not from a
previous version of this code.}
\label{tab:analytic}
\small
\begin{tabular}{llll}
\toprule
Reference & Quantity & Ours & Agreement\\
\midrule
\citet{Wertheim1963}, \citet{Thiele1963} & $c(r)$, PY hard spheres & numerical & first order in $\Delta r$\\
                               & $g(\sigma^+)$, $\eta=0.1$--$0.4$ & numerical & $0.1\,\%$\\
                               & $S(0)=(1-\eta)^4/(1+2\eta)^2$ & numerical & converges\\
Analytic PY, $\eta=0.3$        & $\chi^{-1}$ & 11.07 & 10.66 (extrapolation)\\
                               & $\beta P/\rho$ & 3.890 & 3.816 (grid)\\
                               & PY inconsistency & $+1.33$ & $+1.25$\\
D'Aguanno \& Klein (1991)      & $1/S(0)$, excess pressure & --- & $0.01\,\%$, $0.02\,\%$\\
                               & excess energy, $g_{\max}$ & --- & $0.27\,\%$, $0.54\,\%$\\
\citet{Vrij1979} mixture PY         & $S^{\mathrm{AL}}_{ij}$, 3 classes & numerical & $0.09\,\%$\\
\bottomrule
\end{tabular}
\end{table}

\subsection{Comparison with independent software}

\begin{table}[htbp]
\centering
\caption{Comparison with independently written packages. ``exact'' denotes
agreement to all printed digits; the remaining entries are limited by
discretisation and converge as $\Delta r$ is reduced.}
\label{tab:software}
\small
\begin{tabular}{lllc}
\toprule
Package & Model & Method & Agreement\\
\midrule
\code{mixscatter} \citep{DiazMaier2024} & mixture PY \citep{Vrij1979} & analytic & $0.09\,\%$\\
\code{jscatter} & one-component PY & analytic & $0.02$--$0.12\,\%$\\
\code{sasmodels}, \code{jscatter} & RMSA \citep{HayterPenfold1981} & analytic & $10^{-12}$ (6/6)\\
\code{sasmodels} & two-Yukawa MSA \citep{Liu2005} & analytic & $0.21\,\%$\\
\code{jscatter} & adhesive hard spheres & analytic & see \S\ref{sec:welltrap}\\
\code{jscatter} & OZ solver, PY and HNC & numerical & $0.2$--$0.4\,\%$\\
\code{OrnsteinZernike.jl} & 4 bridge functions & analytic & $\sim10^{-16}$\\
\bottomrule
\end{tabular}
\end{table}

\paragraph{An absolute check: the compressibility.} Every comparison above
sets one calculation against another. One quantity admits an exact answer
instead. For monodisperse hard spheres the Percus--Yevick compressibility is
\begin{equation}
  S(0) = \frac{(1-\varphi)^{4}}{(1+2\varphi)^{2}},
  \label{eq:pyS0}
\end{equation}
so extrapolating the computed partials to the origin and comparing measures
absolute accuracy rather than agreement.

The result is that the transform type matters far more than its place in a
cost table suggests. At $\varphi = 0.40$ and 100 radial points per diameter,
the first-order DST-I gives $S(0)$ low by $5.27\,\%$; the second-order DST-IV
is low by $0.036\,\%$ on the same grid, and by $0.004\,\%$ at four times the
resolution. Type 1 refines as $3.96$ per fourfold step, textbook first order,
so it reaches $1.33\,\%$ where type 4 began 150 times better.

The cause is where the discontinuity lands. A DST-I places nodes at
$r = (n+1)\Delta r$, so the hard core at $r = \sigma$ falls ON a node and is
carried by a point that is neither inside nor outside; the excluded volume is
then wrong at first order in $\Delta r$. A DST-IV places them at
$(n+\frac{1}{2})\Delta r$, the core falls between points, and the step is
resolved to second order.

Two things make this worth reporting rather than absorbing into a
recommendation. First, the error is independent of the real-space range:
extending $r_{\max}$ from $41\sigma$ to $655\sigma$ left $S(0)$ unchanged in
the sixth decimal. A convergence study that refines resolution sees orderly
first-order behaviour and concludes the method works --- which it does, slowly,
towards an answer whose leading error is set by grid alignment. Second,
$S(0)$ is the isothermal compressibility, so this is a systematic bias in
exactly the quantity the volume fraction controls, and a fit will move
$\varphi$ to absorb it.

It also accounts for a discrepancy that had resisted explanation: a
multicomponent solve converged to a fixed point at $7\times10^{-12}$ yet
differing from two independent analytic references by $\sim10^{-3}$, the
difference immune to refinement and growing with $\varphi$. That is the same
effect seen from the other side.

The practical consequence is a restriction rather than a fix. Type 4 needs
every pair core between grid points, which one size class can be given and a
mixture cannot: the pair diameters $\sigma_{ij} = (\sigma_i+\sigma_j)/2$ come
from a moment-matched quadrature and are mutually incommensurate. The
polydisperse route is therefore first order, and until a mixture-capable
second-order treatment exists a dense polydisperse fit should use the finest
affordable resolution, with its fitted volume fraction read accordingly.

\paragraph{Reproducibility of the figures and of the table.} Three gaps are
recorded here rather than left to be discovered.

The remaining \code{jscatter} rows above --- one-component PY, adhesive
hard spheres and the numerical OZ comparison --- were measured where
\code{jscatter} was installed. Its dependency chain does not build under
MinGW, so they cannot be re-run on the machine this work is written on. Two
rows have since been brought back within reach, and by the same move:
replacing a reference that has to be obtained with one that can be
installed.

\paragraph{Two reference solvers replaced by \code{sasmodels}.} The RMSA
and two-Yukawa rows are now measured against SasView's \code{sasmodels},
which is 3-clause BSD and installs with \code{pip} alone. Both previously
depended on something a reader could not simply obtain --- \code{jscatter},
which will not build here, and a 97\,kB vendored copy of the two-Yukawa
solver used under a citation-ware grant the authors gave in writing. A
licence governs an implementation rather than a method, so \citet{Liu2005}
is cited exactly as before; what changes is that anyone can now repeat the
comparison.

Agreement is to $2\times10^{-8}$, $7\times10^{-8}$ and $3\times10^{-8}$ on
the three two-Yukawa state points where both solve cleanly. Of the
remaining two, \code{sasmodels} returns a physical $S(q)$ at
$\varphi = 0.10$ where the vendored solver reports no root found --- a
root-finder limitation rather than a statement about the physics, since the
same parameters at HIGHER volume fraction solve in both. The last,
$Z_1 = 6$ against $Z_2 = 4$, agrees only to $2\times10^{-3}$, with the
discrepancy peaking at the structure-factor maximum; the two inverse ranges
are close there, which is where the quartic's roots crowd together and two
implementations may select marginally differently. Both results are
physical and both are defensible.

\paragraph{The sign of $K$, and where the inconsistency actually was.}
\citet{Liu2005} write the tail as $-K\exp[-Z(r-1)]/r$, so positive $K$ is
ATTRACTIVE. This package's solver follows that convention, and so does
\code{sasmodels}. The vendored reference file did not: reproducing its
output required flipping both signs, after which the two agree to
$10^{-8}$, where every other candidate (length unit, pair ordering, $Z$
scaling) stayed above 20 per cent at every state point.

So the inconsistency was in the comparison apparatus rather than in the
solver being compared --- which is the more flattering of the two
possibilities, and was not the one assumed while the discrepancy was being
chased. Retiring the vendored file removes it entirely, and is a further
reason to prefer \code{sasmodels}: one fewer convention in the tree.

\code{fig3\_oscillation.pdf} has no generator in the repository. The figure
exists and nothing can reproduce it, which means a referee asking for any
change to it would find there is nothing to change.
\code{fig\_survey.pdf} names \code{tools/section53\_survey.py} in its own
caption and is the model the others should follow;
\code{fig1\_discretisation.pdf} now has
\code{tools/fig1\_discretisation.py}.

Regenerating it also caught a discrepancy in the figure itself.
\code{polydisperse\_nodes} offers two discretisations --- \code{sizeClasses},
the moment-matched generalised Gauss--Laguerre rule the caption describes,
and \code{quantileClasses}, which is quantile-based --- and only the first
reproduces the moments exactly. At $p = 3$ the Gauss--Laguerre rule returns
$\langle\sigma\rangle = 1.000000$ and $\langle\sigma^{2}\rangle = 1.062500$,
exactly $1 + s^{2}$; the quantile rule returns $0.9794$ and $0.9596$. Both
are usable discretisations, but only the first supports the claim printed
beside the figure. The moments are now shown on each panel so that the claim
can be checked rather than taken on trust.

\paragraph{The table's figures depend on the transform, and the table does
not say so.} The one-component Percus--Yevick row is quoted as
$0.02$--$0.12\,\%$ over $0.2 \le \varphi \le 0.4$, and that is a DST-IV
result. With the first-order DST-I at the same 100 radial points per
diameter the same comparison gives $0.70$, $1.37$ and $2.33\,\%$ at
$\varphi = 0.2$, $0.3$ and $0.4$ --- an order of magnitude worse, rising with
density, and reducible only fourfold per fourfold refinement
($0.17$, $0.33$, $0.51\,\%$ at 400 points).

This is the grid-alignment effect described above --- the hard core falling ON
a DST-I node rather than between DST-IV nodes --- seen in a second
quantity, and it carries the same caveat: type 4 is available for ONE size
class and not for a mixture, so the polydisperse route --- which is what this
paper is about --- carries the first-order figures rather than the quoted
ones. A reader comparing packages should know which transform produced a
number before comparing it with their own.

\paragraph{RMSA, and a units error that imitated a physics disagreement.}
\label{sec:rmsarescale}
This row agrees with \emph{two} independent implementations to $10^{-12}$
over $\gamma = 3.7$ to $816$ and $\kappa\sigma = 1.5$ to $3.8$, including the
strongly coupled points at which the Hayter--Penfold rescaling engages.
Reaching that took an afternoon, and the detour is worth recording because
the symptom was convincingly misleading.

The earlier claim was ``exact (9/9)'' against \code{jscatter}, which does not
build under MinGW: unrepeatable on the machine this is developed on.
\code{sasmodels}' \code{hayter\_msa} is a better reference --- a different
code lineage, installable with \code{pip} and nothing else --- but comparing
against it gave disagreements of 6 to 32 per cent that GREW WITH COUPLING.
That is the signature of a physics difference rather than a units slip, and
it was pursued as one: root selection between the two solutions the closure
admits, the rescaling threshold, radius-versus-diameter in $\gamma$. Each
was plausible, each was testable, each was wrong. No single rescaling of any
input reconciled the two, and the best-fitting factor itself drifted with
coupling.

The answer is in Eq.~(3a) of \citet{HayterPenfold1981}. The reduced pair
potential is written
\begin{equation}
\beta U(x) = \gamma\,\frac{e^{-kx}}{x},
\qquad x = r/\sigma,\quad k = \kappa\sigma,
\label{eq:hpgamma}
\end{equation}
so the value AT CONTACT is $\gamma e^{-k}$, and that --- not $\gamma$ --- is
what the paper equates to $\beta z^2/[\pi\epsilon_0\epsilon\sigma(2+k)^2]$.
\code{sasmodels} computes the contact potential; \code{jscatter} and the
wrapper here take the amplitude $\gamma$. The two differ by $e^{k}$, and
because $k$ varies with the screening the discrepancy grows with coupling
--- which is exactly why it read as physics.

Two things are worth taking from this. A discrepancy that scales with a
physical parameter is NOT evidence that the difference is physical: a
conversion factor depending on that parameter behaves identically. And the
arbiter was the original paper rather than either code's documentation ---
both describe $\gamma$ as ``the contact potential'', and both use the phrase
loosely for a quantity their own sources define differently.

With the factor applied, \code{tools/rmsa\_vs\_sasmodels.py} reproduces
\code{sasmodels} at all six state points, and the same inputs reproduce
\code{jscatter} to the same precision. The claim in the table is therefore
checkable by any reader with \code{pip install sasmodels}, which the one it
replaces was not.

\paragraph{A closure whose domain is unmapped.} The Carbajal--Tinoco bridge
function \citep{CarbajalTinoco2022} is implicit, $b = T(\gamma + b)$, and
has no solution at all below a
threshold: substituting $w = \gamma + b$ turns it into the explicit map
$\gamma(w) = w - T(w)$, whose range is bounded below. Numerically
$\gamma_{\min} = -0.509$ at $\lambda = 0$, rising towards zero as $\lambda$
grows, with two solutions above it and none below. Neither the source paper
nor the comparison of \citet{PihlajamaaJanssen2024} locates that boundary;
the latter simply excludes closures from its figures ``due to convergence
issues at this state point''.

\subsection{Analytic references implemented in this work}
\label{sec:analyticrefs}

Two of the three classical solvable polydisperse models had no reference
implementation available to us, so we implemented them. Together with the
Vrij hard-sphere solution supplied by \code{mixscatter}, they give exact
analytic coverage of the polydisperse hard-sphere, charged-hard-sphere and
sticky-hard-sphere fluids, against which the numerical solver can be checked
at any state point.

\paragraph{Polydisperse charged hard spheres.}
Appendix A of Gazzillo, Giacometti \& Carsughi (1999), the corrected form of
their 1997 paper --- the correction supplies a factor $\pi$ omitted from
$P_z$ and repairs an error in the determinant. This is the same physics
implemented by the FLAC program of Carsughi \emph{et al.}, which cites the
uncorrected 1997 paper. Our implementation reproduces the neutral
monodisperse limit against analytic Percus--Yevick to $3\times10^{-14}$, the
neutral polydisperse limit against \code{mixscatter} to $1\times10^{-13}$,
and, for charged systems, all three effects the authors report for increasing
polydispersity: the first peak is lowered ($1.7350\to1.3546$), shifted to
smaller $Q$ ($5.841\to5.030$ in units of $Q\langle\sigma\rangle$) and
$S(0)$ increased ($0.0187\to0.0484$), reproducing their component counts (79,
112, 149) exactly.

One limitation of that route is worth stating, because it is a modelling
assumption rather than an approximation that improves with effort. The
corresponding-states construction takes the valence to be PROPORTIONAL TO
SURFACE AREA, so charge polydispersity is fully correlated with size
polydispersity: one distribution governs both. A sample whose surface
chemistry does not scale with area --- differing grafting densities across a
size fraction, say --- cannot be represented.

\citet{GinozaYasutomi1998} treat size and interaction polydispersity as
\emph{independent}, on the same factorisable-coefficient mean-spherical
solution, and would be the route to take if that case ever arose. Both
descend from \citet{Ginoza1986}, who showed that factorisable coefficients
$K_{ij} = K d_i d_j$ collapse the multicomponent Blum--H\o ye equations to
rational functions of a single scalar $\Gamma$ --- which is what makes any of
this affordable for a size distribution, since each component then enters
only through $(\sigma_i, \rho_i, d_i)$ and the cost per evaluation is linear
in the number of classes rather than quadratic.

A partial implementation of the Ginoza route exists in
\code{ginoza\_yukawa\_msa.py} and is deliberately shelved: it solves for
$\Gamma$ and stops short of $S(q)$. Finishing it would reproduce, from
equations transcribed out of a scanned 1986 paper, a structure factor the
Gazzillo route already computes and which is validated above at machine
precision. The two closing conditions Ginoza gives, his Eqs.~(38a) and
(38b), are both implemented there and agree to $10^{-12}$, so the
transcription is at least self-consistent; that is recorded in case the
decoupled-charge case ever makes the module worth completing.

\paragraph{Polydisperse sticky hard spheres.}
\citet{GazzilloGiacometti2000}, a mean-spherical solution in the Baxter sticky
limit with factorisable coupling parameters $K_{ij}=Y_iY_j$, which renders
the Baxter factor matrix three-dyadic and collapses the $p\times p$ inversion
to a $4\times4$ determinant for any number of components. Validation:
the $\gamma_0=0$ limits against analytic Percus--Yevick ($4\times10^{-14}$)
and \code{mixscatter} ($1\times10^{-13}$); the monodisperse sticky
compressibility against their closed form ($3\times10^{-8}$); the critical
point $\eta_c=(\sqrt3-1)/2$, $\gamma_{0c}^2=(\sqrt3+2)/6$ located as the
divergence of $S(0)$; and the generalised Boyle temperatures of their
Eq.~(53) to the printed digits ($2.787$, $1.677$, $2.986$, $2.897$ against
$2.79$, $1.68$, $2.99$, $2.90$).

Note that the stickiness parameter of this model is \emph{not} Baxter's
$\tau$ \citep{Baxter1968}: Baxter's Percus--Yevick solution for mixtures requires $p(p+1)/2$
coupled quadratic equations and is not in general dyadic, which is why the
mean-spherical model exists. The two agree only in the hard-sphere limit.

Three entries in Table~\ref{tab:software} deserve comment.

The \textbf{RMSA} comparison was exact at all nine state points tested
($\eta = 0.05$--$0.25$, screening lengths $2$--$10$\,nm, contact potentials
$20$--$100\,k_BT$), including four at which the Hayter--Penfold rescaling
\citep{HayterPenfold1981}
procedure engaged, so the agreement is not merely that of the underlying MSA
branch. Two independently written implementations of a numerically delicate
algorithm reproduce each other to all printed digits.

The \textbf{two-Yukawa} comparison validates the hard-core double-Yukawa
potential, including cases of mixed sign (attractive $K_1$, repulsive $K_2$),
which is the regime in which a sign convention is most easily mistaken. The
reference implementation states the identical convention,
$V(r) = -k_BT[K_1 e^{-Z_1(r-1)} + K_2 e^{-Z_2(r-1)}]/r$ for $r>\sigma$.

The \textbf{numerical} comparison is the only external check available for the
hypernetted-chain closure, which has no analytic solution. Both solvers are
discretised and differ in transform, mixing scheme and hard-core
representation, so the $0.3\,\%$ agreement is the expected order.

\subsection{Two defects found by external comparison}

\paragraph{Identical hard cores in mixtures.}
Comparison with \code{mixscatter} revealed a discrepancy of $5.2\,\%$ that,
diagnostically, did \emph{not} decrease under grid refinement and was
\emph{linear in the size spread}, vanishing exactly in the monodisperse limit.
The cause was that the hard-sphere potential routine positioned its core by
array index rather than against the radial grid, so it did not see the
rescaled grid on which the polydisperse construction of
equation~(\ref{eq:reducedtail}) relies. Every pair therefore received its core
at the same radius: the mixture was being solved as though all particles were
identical. After correction the agreement is $0.09\,\%$ and converges. It was
the only one of twenty-one potentials written in that form.

\subsection{Two numerical traps in fitting}

Both of the following are properties of the fitting problem rather than of
this implementation, and both produce a converged fit with a plausible
$\chi^2$ and silently wrong parameters. We report them because neither is
evident from the literature and both cost us considerable time.

\paragraph{A rank-revealing linear solve can discard a parameter entirely.}
It is natural to split the fit: the scale, the flat background and any
power-law amplitude enter the intensity linearly and can be obtained exactly
by weighted least squares at every iteration, leaving only the shape
parameters to the nonlinear optimiser. The saving is real. The hazard is
that the design columns need not be comparable in magnitude. A model built
from scattering length densities in $\mathrm{cm}^{-2}$ carries the contrast
squared and so has magnitude $\sim10^{20}$, while the constant background
column is unity --- a condition number near $10^{20}$.
\code{numpy.linalg.lstsq}, like any rank-revealing solver, discards singular
values below $\max(M,N)\epsilon$, so the background and power-law
coefficients are returned as \emph{exactly} zero and the solve reports
success. The symptom is a fit that a hand-chosen background beats by tens of
percent in $\chi^2$, which is indistinguishable from a broken optimiser.
Equilibrating the columns to unit norm before the solve and rescaling
afterwards recovers all coefficients to five figures and changes nothing
where the problem was already well conditioned \citep{GolubVanLoan2013}.

The same scaling defeats the nonlinear optimiser: a parameter of order
$10^{10}$ satisfies \code{least\_squares}'s gradient test after a single
evaluation, so it never moves and is reported with a meaningless
uncertainty. Since the absolute scattering length densities are in any case
perfectly correlated with the scale, we fit their RATIO, which is of order
unity and whose sign is the quantity the data actually determine.

\paragraph{Gaussian resolution smearing fails where it is needed most.} A
small-angle instrument broadens $\mathbf{Q}$ in two dimensions in the
detector plane, and the data are then radially averaged; the distribution of
the magnitude about a true $Q_0$ is Rician, Eq.~\eqref{eq:rician} below, not
Gaussian \citep{Pedersen1990,BarkerPedersen1995}. A Gaussian in $Q$ is its
$Q_0/\sigma \gg 1$ limit and is entirely adequate wherever
$\mathrm{d}Q/Q$ is small.

\begin{equation}
  P(Q\,|\,Q_0) \;=\; \frac{Q}{\sigma^{2}}
  \exp\!\left(-\frac{Q^{2}+Q_0^{2}}{2\sigma^{2}}\right)
  I_0\!\left(\frac{QQ_0}{\sigma^{2}}\right)
  \label{eq:rician}
\end{equation}

It is not always small. For a spherical-shell dataset used to develop this
work, $\mathrm{d}Q/Q$ reaches $0.55$ at the lowest measured point, where a
Gaussian kernel places $44\,\%$ of that point's weight below $Q_{\min}$ and
some of it at \emph{negative} $Q$. The leading factor of $Q$ in
\eqref{eq:rician} suppresses the origin and excludes negative $Q$ by
construction, reducing that to $34\,\%$. Both forms reproduce the unsmeared
curve to $2\times10^{-16}$ as $\mathrm{d}Q \to 0$, which is the check that
catches a mis-normalised kernel. Evaluate \eqref{eq:rician} through the
exponentially scaled $\tilde{I}_0(x) = e^{-|x|}I_0(x)$: the argument reaches
several hundred at these widths, where $I_0$ overflows a double long before
the prefactor cancels it.

A related trap follows. A background term $A\,Q^{-4+d}$ diverges at the
origin while the smearing kernel reaches below the measured range, so the
smeared value at the lowest point is governed by whatever guard prevents
division by zero --- in our case $6.6\times10^{5}$ times the unsmeared value,
and varying by five orders of magnitude as that guard was moved. The power
law describes the measured range and says nothing about $Q \to 0$, where any
real curve turns over; holding it constant below the lowest measured $Q$
removes the arbitrariness.

\begin{quote}\small\itshape
Draft note on citations. Now placed: \code{OrnsteinZernike1914} and
\code{PercusYevick1958} (the equation and closure, \S2),
\code{HansenMcDonald2013}, \code{Anderson1965}, \code{WalkerNi2011},
\code{HayterPenfold1981}, \code{Liu2005}, \code{DiazMaier2024}
(\code{mixscatter}), \code{BarkerPedersen1995} and
\code{GolubVanLoan2013}.

Still in \code{ozgui.bib} and uncited: \code{DuhHenderson1996},
\code{Gazzillo1999}, \code{Gazzillo2006}, \code{KalyuzhnyiCummings2006},
\code{Carsughi2000}. Each has a natural home --- the first in the discussion
of mixing rules, the Gazzillo papers where stickiness and its size
dependence are treated --- but placing a citation where it is not actually
used is worse than leaving it out, so they are left for whoever writes those
sections.

Still absent and worth adding: software citations for NumPy, SciPy, SUNDIALS
and NLopt, all of which the work depends on.

One substantive point from the verification. \citet{BarkerPedersen1995}
conclude that a Gaussian approximation to the resolution is PREFERRED in most
experiments for ease of calculation, and is inadequate only in particular
cases --- of which they name Porod scattering. The argument above should be
framed as identifying one of those known exceptions rather than as a general
correction, which is both more defensible and more useful.
\end{quote}

\subsection{What these figures measure}

A note on what these figures measure. Every agreement quoted in this section
was obtained with the type-1 discrete sine transform, whose grid points fall
\emph{on} the hard core, and which is therefore first order in $\Delta r$. A
type-4 transform places the grid at $(n+\tfrac12)\Delta r$ so that the core
lies between points --- the midpoint rule --- and is second order. For one
component it is both more accurate and cheaper, and both advantages grow with
resolution:

\begin{center}\small
\begin{tabular}{rrrr}
\toprule
points per $\sigma$ & type-4 time & accuracy gain\\
\midrule
100 & $0.73\times$ & $139\times$\\
400 & $0.65\times$ & $547\times$\\
800 & $0.49\times$ & $1092\times$\\
\bottomrule
\end{tabular}
\end{center}

The times are relative to type 1 at the same points per diameter; the
accuracy gain is the ratio of the two errors. The growth of that ratio with
resolution is the second-order against first-order signature. The figures
quoted above would improve correspondingly.

The two transforms require different grid \emph{lengths}, which is easy to
overlook and costly to get wrong: a DST-I is computed through an FFT of
length $2(N+1)$ and so wants $N = 2^{k}-1$, whereas a DST-IV is a length-$N$
transform and wants $N = 2^{k}$. Sizing both for the former leaves the latter
on a composite length ($16383 = 3\times43\times127$) and costs it a factor of
three in the transform, enough to make the second-order method appear slower
than the first-order one. Accuracy is unaffected by the choice.

We nevertheless retain type 1 as the default, because the type-4 condition
cannot be met for a mixture. It requires every pair core $\sigma_{ij}$ to
fall between grid points, and with size classes drawn from a Gaussian
quadrature the $\sigma_i$ are irrational, so the $p(p+1)/2$ conditions cannot
be satisfied simultaneously by any choice of resolution. Type 4 is available
as an option and is correct for one component, where the single core
coincides with the reduced unit by construction.


The defect had been invisible to every existing test because the resulting
error vanishes in precisely the monodisperse limit where those tests looked.

\paragraph{Under-resolved narrow wells.}
\label{sec:welltrap}
Comparison with an analytic adhesive-sphere solution showed errors that are
governed not by the number of grid points per particle diameter but by the
number \emph{across the attractive well}:

\begin{center}\small
\begin{tabular}{cc}
\toprule
points across well & max.\ relative error\\
\midrule
2 & 0.916\\
4 & 0.324\\
8 & 0.142\\
16 & 0.068\\
32 & 0.035\\
\bottomrule
\end{tabular}
\end{center}

This is ordinary first-order convergence in $\Delta r/\delta$ and not a
defect, but the default discretisation violates it badly: at 100 points per
diameter a well of width $\delta = 0.02\sigma$ spans two grid points and the
structure factor is in error by ninety per cent, without any warning. We
recommend $\;\gtrsim 30\,\sigma/\delta$ points per diameter for any potential
with a narrow feature, and note that this requirement applies equally to
square-well and piecewise-constant potentials.

\subsection{Internal consistency}

The following are not independent checks but constrain the implementation:
reduction of the polydisperse construction to a single size class reproduces
the original one-component routines bit-identically; the extended
Rogers--Young closure reduces to Rogers--Young to $2.5\times10^{-16}$ when its
additional parameter vanishes; the hard-core double-Yukawa potential agrees
with the three-Yukawa form to $\sim10^{-16}$ under the stated
reparameterisation; and the exact intensity reproduces brute-force
integration of $\langle|F|^2\rangle$ in the dilute limit
(ripple $0.041$ against $0.038$).

% ---------------------------------------------------------------------
% References to add to the bibliography
% ---------------------------------------------------------------------
% Gazzillo, D., Giacometti, A. & Carsughi, F. (1997). J. Chem. Phys. 107,
%   10141. [polydisperse charged hard spheres, MSA]
% Gazzillo, D., Giacometti, A. & Carsughi, F. (1999). arXiv:cond-mat/9909153.
%   [corrected Appendix A; scaling vs decoupling approximation]
% Gazzillo, D. & Giacometti, A. (2000). arXiv:cond-mat/0011330.
%   [polydisperse sticky hard spheres, MSA, factorisable coupling]
% Carsughi, F., Giacometti, A. & Gazzillo, D. (2000). arXiv:cond-mat/0011242.
%   [the FLAC program]
% Tutschka, C. & Kahl, G. (1998). J. Chem. Phys. 108, 9498.
% Blum, L. & Hoye, J. S. (1978). J. Stat. Phys. 19, 317.
%   [multicomponent multi-Yukawa MSA; the finite-range generalisation]
% Hoye, J. S. & Stell, G. [SCOZA]; Reiner, A. & Hoye, J. S. (2008).
%   arXiv:0712.3739. [SCOZA near the critical point]
% Bergmann, M. J. et al. mixscatter (Python package).
%   [add version/DOI; used for the analytic Vrij mixture PY]
% Biehl, R. (2019). jscatter. PLoS ONE 14, e0218789.
%   [Two_Yukawa after Liu et al. (2005); RMSA after Hayter & Penfold (1981)]
% Hayter, J. B. & Penfold, J. (1981). Mol. Phys. 42, 109-118.
% Liu, Y., Chen, W.-R. & Chen, S.-H. (2005). J. Chem. Phys. 122, 044507.
%   [two-Yukawa analytic solution]
% Pihlajamaa, I. & Janssen, L. M. C. (2024). Phys. Rev. E 110, 044608.
%   [closure comparison; OrnsteinZernike.jl]
% Vrij, A. (1979). J. Chem. Phys. 71, 3267-3270.
%   [mixture PY, the reference for the mixture comparison]
% Wertheim, M. S. (1963). Phys. Rev. Lett. 10, 321.
% Thiele, E. (1963). J. Chem. Phys. 39, 474.
% Baxter, R. J. (1968). J. Chem. Phys. 49, 2770.
% Malijevsky, A. & Labik, S. (1987). Mol. Phys. 60, 663.
